% !TeX root = ../../Book3.tex
% 纲要标题：## 第20章　正规性与 Serre 判别法
% 纲要位置：计划纲要.md，第 2584 行。

\chapter{正规性与 Serre 判别法}
\label{b3:ch:20}
\BookChapterNumber{b3:ch:20}

\begin{chapterabstract}
正规性要求每个局部环都是整闭整环。Serre 判据把它分成余维一的正则性与各处的深度条件，并适用于含多个不可约分量的 Noether 环。本章证明这一判据，再说明正规整环为何可由高度一处的 DVR 恢复，以及除子类群如何记录局部阶数不能由一个全局函数实现的障碍。尖点、交叉直线与二次锥则说明正规性、正则性和 CM 性有不同要求。
\end{chapterabstract}

\begin{learninggoals}
\begin{itemize}
\item 准确区分 \(R_k\) 与 \(S_k\) 的量词、局部维数及其表达的条件。
\item 用关联素理想、非零因子和全分式环证明一般 Noether 环的 Serre 正规性判据。
\item 从高度一局部赋值检验环元素，定义 Weil 除子并辨认类群的意义。
\item 用尖点、交叉直线与二次锥比较正规、正则与 CM 性，区分正规化和消解奇点。
\end{itemize}
\end{learninggoals}

设 \(k\) 为域。尖点环 \(R=k[t^2,t^3]\) 的分式域里有 \(t=t^3/t^2\)，而 \(t\) 满足系数在 \(R\) 中的首一方程 \(T^2-t^2=0\)，却不属于 \(R\)。这说明整环即使没有零因子，也可能缺少本应由整方程控制的元素。正规化把这些元素补进来；若原环已经包含全部这样的元素，才称得上整闭。

\ProseParagraphBreak

整闭性在 Noether 环中可以从局部信息辨认，但只在高度一处检查并不总够：较高余维中还可能藏有需要深度条件排除的缺口。Serre 判据将这两种检查分开。对正规整环，高度一处的离散赋值给出每个非零有理函数的局部阶数；这些阶数能否同时来自一个全局函数，便是除子类群所量度的障碍。

\ProseParagraphBreak

本章使用\chapref{b3:ch:04}的关联素理想、\chapref{b3:ch:06}的约化环全分式环、\chapref{b3:ch:10}的 DVR，以及\chapref{b3:ch:16}的深度与 Serre 条件。完全交和 CM 例子采用\chapref{b3:ch:19}，正则局部环的唯一分解性采用\cref{b3:thm:regular-local-ufd}。\secref{b3:sec:2005}另用第二卷第4章已证明的主理想整环模分类，以及有限投射模的对偶和张量积。

\ProseParagraphBreak
松村英之的《Commutative Ring Theory》对正规判据的完整证明见\cite[Theorem 23.8, pp. 183--184]{Matsumura1989CommutativeRingTheory}，DVR 的整闭刻画见\cite[Theorem 11.2, pp. 79--80]{Matsumura1989CommutativeRingTheory}，除子类群与 Picard 群的环论解释见\cite[\S20, pp. 165--167]{Matsumura1989CommutativeRingTheory}。本章保留一般非整环版本，同时把证明中清除分母的步骤单列为引理。

% !TeX root = ../../Book3.tex
% 纲要标题：### 20.1 R_k 与 S_k 条件
% 纲要位置：计划纲要.md，第 2586 行。

\section{\texorpdfstring{\(R_k\)}{Rₖ} 与 \texorpdfstring{\(S_k\)}{Sₖ} 条件}
\label{b3:sec:2001}

正规性的局部检验同时涉及低余维处的环结构与较高余维处的深度。先分别固定 \(R_k\) 和 \(S_k\) 的量词，再用算例说明二者不能互相替代。


% 纲要内容（仅作写作计划，尚非教材正文）：
%
% * 余维上的正则性
% * 深度条件的局部化表述
% * 余维一和高余维承担的不同任务；不可约分量在闭点相交的深度不足例子
%

\subsection{余维上的正则性}
\label{b3:subsec:2001:01}

整闭性表面上要求检查分式域中的所有首一方程，正则性则比较局部环的维数与极大理想生成元数。Serre 判据把前一种条件分成两项：低余维处检查正则性，较高余维处检查深度。先把两项条件的量词固定下来。

\begin{definition}\label{b3:def:serre-rk}
设 \(R\) 为交换 Noether 环，\(k\ge0\) 为整数。若对每个满足 \(\operatorname{ht}\mathfrak p\le k\) 的素理想 \(\mathfrak p\)，局部环 \(R_{\mathfrak p}\) 都正则，则称 \(R\) 满足 \term[Serre-R-tiaojian]{Serre 条件 \(R_k\)}[Serre's condition \(R_k\)]。
\end{definition}

这里的高度是 \(\dim R_{\mathfrak p}\)，不是 \(\dim R/\mathfrak p\)。在仿射整簇中，它衡量所对应子簇的余维。条件 \(R_0\) 表示极小素点处的局部环为域；条件 \(R_1\) 还要求每个高度一素点处的局部环为 DVR，因为一维正则局部环的极大理想由一个元素生成。

\ProseParagraphBreak
条件 \(R_0\) 只保证在各不可约分量的一般点处没有幂零元，不能排除集中在特殊点的幂零元。例如对域 \(k\)，
\(k[x,y]/(x^2,xy)\) 在唯一极小素理想 \((x)\) 处把 \(y\) 变成单位，从而 \(x=0\)，局部化为域 \(k(y)\)；但原环中的 \(x\) 非零且平方为零。

\subsection{深度条件的局部化表述}
\label{b3:subsec:2001:02}

回顾\secref{b3:sec:1604}：交换 Noether 环 \(R\) 满足 \(S_k\)，是指对每个素理想 \(\mathfrak p\)，都有
\[
\operatorname{depth}R_{\mathfrak p}\ge\min\{k,\dim R_{\mathfrak p}\}.
\]
特别地，\(S_1\) 排除嵌入关联素理想；\(S_2\) 在所有维数至少二的局部环中要求能连续取两个非零因子。在维数一处，它只要求一个；维数零处没有额外要求。

\begin{proposition}\label{b3:prop:serre-conditions-localize}
设 \(R\) 为交换 Noether 环，\(k\ge0\) 为整数。条件 \(R_k\)、\(S_k\) 均在局部化后保持；它们也都可在各个极大理想处的局部环上检查。若 \(R\) 为 CM 环，则它满足每个 \(S_k\)。
\end{proposition}
\begin{proof}
若 \(S^{-1}\mathfrak p\) 是 \(S^{-1}R\) 的素理想，则
\((S^{-1}R)_{S^{-1}\mathfrak p}\cong R_{\mathfrak p}\)；其以下的素理想链也都避开 \(S\)，所以高度不变。于是局部条件完全相同。每个素理想包含在某个极大理想中，进一步局部化也得到同一个 \(R_{\mathfrak p}\)，证明检查极大理想处足够。CM 环在各处深度等于维数，因此满足所列全部不等式。
\end{proof}

\begin{proposition}\label{b3:prop:reduced-serre-criterion}
对交换 Noether 环 \(R\)，约化性等价于 \(R_0\) 与 \(S_1\) 同时成立。
\end{proposition}
\begin{proof}
若 \(R\) 约化，则其极小素点处为约化零维局部环，因而为域。
\cref{b3:prop:reduced-ring-no-embedded-primes}和\cref{b3:thm:serre-s1-associated}给出 \(S_1\)。

反之，若幂零根 \(N\) 非零，则它作为有限生成的 \(R\) 模有一个关联素理想 \(\mathfrak p\)。因为 \(N\subseteq R\)，有 \(\mathfrak p\in\operatorname{Ass}R\)；\(S_1\) 使它为极小素理想。于是 \(N_{\mathfrak p}\ne0\)，但 \(R_0\) 使 \(R_{\mathfrak p}\) 为域，其幂零根只能为零，矛盾。
\end{proof}

\subsection{余维一与高余维的条件比较}
\label{b3:subsec:2001:03}

\(R_1\) 与 \(S_2\) 不能互相替代。设 \(k\) 为域。一维尖点局部环
\(k[t^2,t^3]_{(t^2,t^3)}\) 是 CM 环，故满足 \(S_2\)；它却不是 DVR，所以不满足 \(R_1\)。另一类例子由两张平面只在一个点相交得到。

\begin{example}\label{b3:exm:r1-not-s2}
设 \(k\) 为域，
\[
T=k[x,y,z,w]_{(x,y,z,w)},\qquad
R=T/\bigl((x,y)\cap(z,w)\bigr).
\]
则 \(R\) 为二维约化局部环，满足 \(R_1\)，却有 \(\operatorname{depth}R=1\)，因而不满足 \(S_2\)。
\end{example}
\begin{proof}
\(R\) 的极小素理想为 \((x,y)R\) 与 \((z,w)R\)。它们的商均为二维正则局部环，且两理想之和为 \(R\) 的极大理想，故两个不可约分量只在闭点相遇，\(\dim R=2\)。设闭点为 \(\mathfrak m\)。若 \(\mathfrak p\ne\mathfrak m\)，至少有一个坐标在 \(\mathfrak p\) 外；把它逆转后，相对的两个坐标由交理想中的乘积关系变为零，局部环成为上述一个正则局部环的局部化。
\cref{b3:thm:regular-local-localization}使它正则，故 \(R_1\) 成立。

自然差映射给出正合列
\[
0\longrightarrow R\longrightarrow T/(x,y)\oplus T/(z,w)
\xrightarrow{(a,b)\mapsto a-b}k\longrightarrow0.
\]
中间模上 \(x+z,y+w\) 依次为非零因子，在两个直和分量中分别是两个剩余变量，所以它的深度为二；右端的深度为零。\cref{b3:thm:depth-lemma}给出 \(\operatorname{depth}R=1\)：深度至少一来自中间模深度，若至少二，则对上述正合列应用 \(\operatorname{Hom}_R(k,-)\) 的长正合列会使从非零的 \(\operatorname{Hom}_R(k,k)\) 到零的 \(\operatorname{Ext}_R^1(k,R)\) 为单射，矛盾。
\end{proof}

这个深度计算也记录了函数的延拓障碍。令 \(U=\operatorname{Spec}R\setminus\{\mathfrak m\}\)。删去交点后，两个分量互不相交，且各自在 \(U\) 中既开又闭；在第一个分量上取常数零、第二个分量上取常数一，便得到 \(U\) 上的正则函数。它不能由 \(R\) 中的元素给出：上面正合列将 \(R\) 识别为两份局部平面环中在闭点剩余值相等的函数对，而这两个常数的差为 \(-1\in k\)，并不为零。

\ProseParagraphBreak
余维一检查没有看见只发生在闭点的相交；这里不满足 \(S_2\)，与这项高余维处的不相容函数障碍相对应。Serre 判据将证明，对于整闭性，余维一的正则性与 \(S_2\) 合起来已经足够。

% !TeX root = ../../Book3.tex
% 纲要标题：### 20.2 Serre 正规性判据
% 纲要位置：计划纲要.md，第 2592 行。

\section{Serre 正规性判据}
\label{b3:sec:2002}

分式在各高度一局部环中属于原环，还需要深度条件才能返回全环。Serre 判据把这一返回过程与一维局部环的 DVR 结构结合，给出适用于具有多个不可约分量的 Noether 环的正规性判别。


% 纲要内容（仅作写作计划，尚非教材正文）：
%
% * Noether 环上 R_1 与 S_2 的组合
% * 约化、全分式环整闭、局部环为整环的三步证明；以坐标幂等元排除不可约分量的局部相交
% * 正规局部整环的深度二与行列式技巧；汇总 Serre 条件、约化性与正规性的对应
%

\subsection{Noether 环上的 \texorpdfstring{\(R_1\)}{R1} 与 \texorpdfstring{\(S_2\)}{S2} 条件}
\label{b3:subsec:2002:01}

\begin{definition}\label{b3:def:normal-noether-ring}
设 \(R\) 为交换 Noether 环。若对每个素理想 \(\mathfrak p\)，\(R_{\mathfrak p}\) 都是在自身分式域中整闭的整环，则称 \(R\) 为\term[zhengguihuan]{正规环}[normal ring]。当 \(R\) 为整环时，这与 \(R\) 在 \(\operatorname{Frac}(R)\) 中整闭等价。
\end{definition}

本书的“正规环”采用这项逐点约定，允许原环有多个不可约分量；整环情形与分式域整闭性的等价已由\cref{b3:prop:normality-local-detection}证明。一般正规 Noether 环的不同不可约分量必须互不相交，而且各自是开闭子集，见下面的\cref{b3:cor:normal-product-domains}。例如，设 \(k\) 为域，则 \(k\times k\) 正规，而曲线 \(xy=0\) 在原点的两条局部分支相交，局部环 \(k[x,y]_{(x,y)}/(xy)\) 不正规。本节同时处理这种非整情形。

\ProseParagraphBreak
\(S_2\) 在判据中要解决的是清除分母。若 \(a\) 是非零因子，而 \(b\notin aR\)，非零商类 \(\bar b\in R/aR\) 必在某个关联素理想处仍然非零。下面两个引理说明，\(S_2\) 迫使这样的素理想高度为一。因此，一个分母无法清除的障碍不会只在高度至少二的素点处出现。

\begin{lemma}\label{b3:lem:s2-principal-associated-height}
设 \(R\) 为满足 \(S_2\) 的交换 Noether 环，\(a\in R\) 为非零因子且非单位。则
\[
\mathfrak p\in\operatorname{Ass}_R(R/aR)
\quad\Longrightarrow\quad\operatorname{ht}\mathfrak p=1.
\]
\end{lemma}
\begin{proof}
由\cref{b3:prop:associated-localization}，\(\mathfrak pR_{\mathfrak p}\) 是 \((R/aR)_{\mathfrak p}\) 的关联素理想，\cref{b3:cor:depth-zero-socle}给出这个非零局部模的深度为零。因为 \(a\in\mathfrak p\)，它在 \(R_{\mathfrak p}\) 上仍是非零因子且非单位。\cref{b3:cor:regular-quotient-depth}说明商去这个正则元素使深度降一，因而
\[
\operatorname{depth}R_{\mathfrak p}
=\operatorname{depth}_{R_{\mathfrak p}}(R/aR)_{\mathfrak p}+1=1.
\]
条件 \(S_2\) 要求 \(1\ge\min\{2,\dim R_{\mathfrak p}\}\)，故 \(\dim R_{\mathfrak p}\le1\)；深度不超过维数又使其至少为一。于是 \(\operatorname{ht}\mathfrak p=\dim R_{\mathfrak p}=1\)。
\end{proof}

\begin{lemma}\label{b3:lem:s2-fraction-membership}
设 \(R\) 为满足 \(S_2\) 的交换 Noether 环，\(a\in R\) 为非零因子，\(b\in R\)。若对每个高度一素理想 \(\mathfrak p\) 都有 \(b\in aR_{\mathfrak p}\)，则 \(b\in aR\)。
\end{lemma}
\begin{proof}
若 \(a\) 为单位，结论直接成立。否则反设 \(b\notin aR\)，则商类 \(\bar b\) 生成 \(R/aR\) 的一个非零有限生成子模，存在
\(\mathfrak p\in\operatorname{Ass}(R\bar b)\subseteq\operatorname{Ass}(R/aR)\)。上引理使 \(\operatorname{ht}\mathfrak p=1\)。在该处 \((R\bar b)_{\mathfrak p}\ne0\)，而它由 \(\bar b/1\) 生成，所以 \(\bar b/1\ne0\)。这正是 \(b/1\notin aR_{\mathfrak p}\)，与假设相矛盾。
\end{proof}

这一步把“某个分式不是环元素”变成“在某个高度一素点已经不是环元素”。关键是利用非零商类所生成子模的关联素理想；\(S_2\) 保证这个关联素理想恰好具有所需高度。\(R_1\) 随后提供这些高度一局部环的 DVR 结构，两项条件由此分别承担检测分母与清除分母的任务。

\subsection{局部环整闭性与 Serre 判据}
\label{b3:subsec:2002:02}

要从正规性推出 \(S_2\)，一维局部整环中的第一个非零因子很容易取得，困难在于局部维数至少二时再选第二个。商去第一个元素后，若极大理想成了关联素理想，就会得到一个不属于原环的分式。证明将按这个分式是否把极大理想映入自身分成两种情形：前者违背整闭性，后者会使极大理想成为主理想，违背维数假设。

\begin{lemma}\label{b3:lem:normal-local-depth-two}
设 \((R,\mathfrak m)\) 为交换 Noether 整闭局部整环。若 \(\dim R\ge2\)，则 \(\operatorname{depth}R\ge2\)。
\end{lemma}
\begin{proof}
选 \(0\ne a\in\mathfrak m\)，它是非零因子。只需证明
\(\mathfrak m\notin\operatorname{Ass}(R/aR)\)，然后用\cref{b3:lem:finite-prime-avoidance}在 \(\mathfrak m\) 中选第二个非零因子。

若反而存在 \(b\notin aR\) 使 \(\mathfrak mb\subseteq aR\)，置 \(z=b/a\in\operatorname{Frac}(R)\setminus R\)，则 \(z\mathfrak m\subseteq R\)。由于 \(\dim R\ge2\)，极大理想 \(\mathfrak m\ne0\)；又因 \(R\) 是整环，\(\operatorname{Ann}_R(\mathfrak m)=0\)，所以 \(\mathfrak m\) 是有限生成忠实模。假如 \(z\mathfrak m\subseteq\mathfrak m\)，乘以 \(z\) 就给出 \(\mathfrak m\) 的一个 \(R\) 线性自同态。这正使行列式技巧适用：一个分式稳定作用于这里的有限生成忠实理想时，其乘法自同态被首一多项式零化。对\cref{b3:thm:determinant-trick}取其中的理想 \(I=R\)，得到首一多项式 \(P(T)\in R[T]\)，使
\[
P(z)\mathfrak m=0.
\]
选取 \(0\ne v\in\mathfrak m\)，在分式域中由 \(P(z)v=0\) 得 \(P(z)=0\)。因此 \(z\) 对 \(R\) 整，整闭性使 \(z\in R\)，矛盾。

若 \(z\mathfrak m\) 不包含于 \(\mathfrak m\)，由于仍有 \(z\mathfrak m\subseteq R\)，便有 \(u\in\mathfrak m\) 使 \(zu\in R\setminus\mathfrak m\)，即 \(zu\) 为单位。对任意 \(v\in\mathfrak m\)，
\[
\frac vu=\frac{zv}{zu}\in R,
\]
故 \(\mathfrak m=uR\)。\cref{b3:thm:principal-ideal-theorem}使 \(\dim R\le1\)，又矛盾。因此 \(R/aR\) 的极大理想不是关联素理想，存在 \(c\in\mathfrak m\) 在 \(R/aR\) 上正则，\(a,c\) 即为长度二的正则序列。
\end{proof}

证明反向蕴含时，需要分别得到约化性、全分式环中的整闭性和每个局部环的整环性。\(R_0\) 与 \(S_1\) 先排除幂零元，高度一处的整闭性和 \(S_2\) 再把整分式追回原环。此时还没有得到局部环是整环；最后须用全分式环的坐标幂等元，把不同分量的分离带回原环，再利用局部性排除多个分量同时存在。

\begin{theorem}[Serre 正规性判据]\label{b3:thm:serre-normality}
设 \(R\) 为交换 Noether 环，不要求 \(R\) 为整环。则
\[
R\;\text{正规}\quad\Longleftrightarrow\quad R\;\text{同时满足}\;R_1\;\text{与}\;S_2.
\]
\end{theorem}
\begin{proof}
零环时两侧都是空条件；以下设 \(R\ne0\)。正规性由各素点的局部环定义，\(R_1\) 与 \(S_2\) 又由\cref{b3:prop:serre-conditions-localize}在局部化后保持，故可在任意给定素点处再局部化。

先设 \(R\) 正规。在高度零处局部环为域；在高度一处，一维 Noether 整闭局部整环由\cref{b3:thm:dvr-characterizations}为 DVR，所以 \(R_1\) 成立。对任意素点，若局部维数为一，整环中可取非零非单位，深度至少一；若局部维数至少二，上引理给出深度至少二；维数零处无要求。这证明 \(S_2\)。

反之，设 \(R_1\) 与 \(S_2\) 成立。它们分别蕴含 \(R_0\) 与 \(S_1\)，因此\cref{b3:prop:reduced-serre-criterion}给出 \(R\) 约化。

接着检验全分式环中的整闭性。取 \(Q(R)\) 中一个对 \(R\) 整的元素 \(z=b/a\)，其中 \(a\) 为非零因子。在每个高度一素理想 \(\mathfrak p\) 处，\(R_1\) 使 \(R_{\mathfrak p}\) 成为一维正则局部环。它由\cref{b3:prop:regular-local-domain}为整环，极大理想又由一个元素生成，故\cref{b3:thm:dvr-characterizations}使它成为 DVR，因而整闭。\(a/1\) 在这里仍为非零因子，分式 \(b/a\) 的像属于其分式域，并且仍满足同一个首一方程，所以属于 \(R_{\mathfrak p}\)。现在已在每个高度一处清除了分母，\cref{b3:lem:s2-fraction-membership}给出 \(b\in aR\)，即 \(z\in R\)。这证明 \(R\) 在 \(Q(R)\) 中整闭。

最后还须从全分式环的整闭性得到局部环是整环。把 \(R\) 局部化后，上述论证仍适用，因而现在可设 \(R\) 局部。\cref{b3:prop:reduced-total-quotient-product}给出
\[
Q(R)\cong\prod_{\mathfrak p\in\operatorname{Min}R}\operatorname{Frac}(R/\mathfrak p).
\]
若这个乘积有至少两个因子，则任一坐标幂等元 \(e_i=(0,\ldots,1,\ldots,0)\) 都不同于 \(0,1\)。它满足 \(e_i^2-e_i=0\)，对 \(R\) 整，刚证明的整闭性迫使 \(e_i\in R\)。但局部环只有幂等元 \(0,1\)：\(e_i+(1-e_i)=1\) 保证两者至少一个为单位，而 \(e_i(1-e_i)=0\) 随即迫使另一个为零，矛盾。因此上面的乘积只能有一个因子。约化性使全部极小素理想的交为零，所以唯一极小素理想就是零，\(R\) 是整环；已经证明的全分式环整闭性便成为分式域整闭性。各点局部环都如此，故原环正规。
\end{proof}

上述\term[Serre-zhenggui-xing-panju]{Serre 正规性判据}[Serre's criterion for normality]可与松村英之的证明比较\cite[Theorem 23.8, pp. 183--184]{Matsumura1989CommutativeRingTheory}。幂等元步骤把全分式环中的整闭性转化为定义所需的局部整环性，说明不同不可约分量在正规环中不能相交。

\subsection{非零因子与深度方法}
\label{b3:subsec:2002:03}

\begin{corollary}\label{b3:cor:normal-product-domains}
非零交换 Noether 环 \(R\) 正规，当且仅当它是有限个 Noether 整闭整环的直积。
\end{corollary}
\begin{proof}
若 \(R\) 正规，定理证明已说明它约化且在全分式环中整闭。全分式环的各坐标幂等元 \(e_1,\ldots,e_s\) 都属于 \(R\)，于是
\(R\cong Re_1\times\cdots\times Re_s\)。每个 \(Re_i\) 是 \(R\) 的商环，故为 Noether 环；它嵌入相应的分式域，所以是整环。若该分式域中的元素 \(\alpha\) 对 \(Re_i\) 整，把它放在第 \(i\) 个坐标、其余坐标置零，得到 \(\tilde\alpha\in Q(R)\)。将其首一方程的各低次系数看作 \(Re_i\subseteq R\) 的元素，最高次系数取 \(1_R\)，同一方程就在每个坐标中零化 \(\tilde\alpha\)。因此 \(\tilde\alpha\) 对 \(R\) 整，属于 \(R\)；又有 \(e_i\tilde\alpha=\tilde\alpha\)，故仍在 \(Re_i\) 中。反向由直积环的每个局部环都是一个因子的局部环，且整闭整环的局部化仍整闭得到。
\end{proof}

这里每个因子 \(Re_i\) 同时对应 \(\operatorname{Spec}R\) 中的 \(D(e_i)=V(1-e_i)\)，所以它既开又闭；因子为整环，又使这个子集不可约。这些开闭不可约分量互不相交并覆盖整个素谱。两张平面在一点相交的例子便违反了这一必要条件，即使它在所有高度一素点处都正则，仍不能正规。

\ProseParagraphBreak
各项局部条件与两个整体性质的关系汇总于\cref{b3:tab:serre-conditions}。

\begin{center}
\begin{minipage}{\textwidth}
\makeatletter
\def\@captype{table}
\makeatother
\centering
\caption[Serre 条件与约化性、正规性]{Serre 条件与约化性、正规性。设 \(R\) 为交换 Noether 环；每一行的左右两项等价，\(\mathfrak p\) 均遍历 \(R\) 的素理想。}
\label{b3:tab:serre-conditions}
\begin{tabularx}{\textwidth}{@{}>{\raggedright\arraybackslash}p{0.22\textwidth}>{\raggedright\arraybackslash}X@{}}
\toprule
条件 & 含义或等价性质\\
\midrule
\(R_0\) & 每个极小素理想处的局部环为域。\\
\(S_1\) & \(\operatorname{Ass}_R R=\operatorname{Min}R\)，即无嵌入关联素理想。\\
\(R_0\) 与 \(S_1\) & \(R\) 约化。\\
\(R_1\) & 高度不超过一处的局部环正则：高度零处为域，高度一处为 DVR。\\
\(S_2\) & 对每个 \(\mathfrak p\)，\(\operatorname{depth}R_{\mathfrak p}\ge\min\{2,\dim R_{\mathfrak p}\}\)。\\
\(R_1\) 与 \(S_2\) & \(R\) 正规。\\
\bottomrule
\end{tabularx}
\end{minipage}
\end{center}

表中 \(R_0\) 与 \(R_1\) 检验局部环的结构，\(S_1\) 与 \(S_2\) 检验深度；约化性和正规性都需要把这两类条件配合起来。

\ProseParagraphBreak
判据中的各个条件因而各有明确作用。Noether 性保证关联素理想有限、商模有可选的关联素理想，并使行列式技巧作用于有限生成极大理想。\(R_1\) 使高度一局部环成为整闭的 DVR。\(S_2\) 则保证分母清不掉的障碍已经出现在高度一处，而不是只藏在较高余维。

\ProseParagraphBreak
这也解释了为什么不能只凭“奇点集合余维至少二”就断言正规：该句话至多提供 \(R_1\)，还必须另有 \(S_2\)。例如 CM 环自动提供 \(S_2\)，于是对 CM 环，正规性确实等价于 \(R_1\)；\cref{b3:exm:r1-not-s2}中两张平面只在一点相交，正是缺少这项深度条件。正则环的所有局部环都正则，因而既满足 \(R_1\)，又由\cref{b3:cor:regular-ring-cm}为 CM 环、满足 \(S_2\)，所以正则环必正规。

% !TeX root = ../../Book3.tex
% 纲要标题：### 20.3 高度一局部环与除子类群
% 纲要位置：计划纲要.md，第 2598 行。

\section{高度一局部环与除子类群}
\label{b3:sec:2003}

正规整环的高度一局部环是 DVR，每个这样的素理想因此给出一个离散阶数。这些阶数可以检测分式是否属于原环，也能记录主理想不能表示的除子信息。


% 纲要内容（仅作写作计划，尚非教材正文）：
%
% * 正规整环是分式域中全部高度一局部环的交；用负赋值检测分式不属于环
% * 归一化离散赋值；正阶对应主理想的极小素理想，故主除子支撑有限
% * Weil 除子作为有限整数阶数资料，类群为零回用高度一素理想主性判据
% * 函数与普通理想的高度一恢复不同；类群局部化为删除素理想类所得的商群
%

\subsection{正规整环与高度一局部化}
\label{b3:subsec:2003:01}

\begin{theorem}\label{b3:thm:normal-height-one-intersection}
设 \(R\) 为交换 Noether 正规整环，\(K=\operatorname{Frac}(R)\)。则在 \(K\) 中有
\[
R=\bigcap_{\operatorname{ht}\mathfrak p=1}R_{\mathfrak p},
\]
其中交集遍历 \(R\) 的所有高度一素理想 \(\mathfrak p\)，每个 \(R_{\mathfrak p}\) 都是 DVR。若没有高度一素理想，则 \(R\) 为域，交集按 \(K\) 解释。
\end{theorem}
\begin{proof}
DVR 结论由正规性及\cref{b3:thm:dvr-characterizations}得到。包含关系从左到右显然；反向把交集中的元素写成 \(b/a\)，其中 \(0\ne a\in R\)。Serre 判据使 \(R\) 满足 \(S_2\)，所以\cref{b3:lem:s2-fraction-membership}使 \(b\in aR\)。若无高度一素理想，任意非零非单位的主理想都应有一个高度一极小素理想，故不存在这样的非单位，\(R\) 是域。
\end{proof}

记 \(v_{\mathfrak p}:K^*\rightarrow\mathbb Z\) 为 \(R_{\mathfrak p}\) 的归一化离散赋值，即统一化元的值为 \(1\)。对 \(0\ne f\in K\)，上述定理等价于
\[
f\in R\quad\Longleftrightarrow\quad v_{\mathfrak p}(f)\ge0
\quad\text{对所有高度一}\;\mathfrak p.
\]
并且 \(f\in R^*\) 当且仅当所有赋值均为零，因为这时 \(f\) 与 \(f^{-1}\) 同时属于 \(R\)。环元素与单位因此可由各个余维一处的阶数判定。从函数的角度看，一个分式若在所有余维一处都没有极点，就已经属于原环，无须再到余维至少二的点作额外检测。正规性使函数的高余维延拓由这些阶数控制。

\ProseParagraphBreak
例如，设 \(k\) 为域，\(R=k[x,y]\)。在 \(R_{(y)}\) 中，\(y\) 是统一化元而 \(x\) 是单位，所以
\[
v_{(y)}(x/y)=0-1=-1.
\]
仅这一个高度一素理想便检测到 \(x/y\notin R\)。\((x,y)\) 的高度为二，虽然它在几何上是一个特殊的闭点，却不是这里给出离散阶数的位置。高余维处的理想信息并未消失；对有理函数的环元素判定，高度一检查已经足够。

\subsection{Weil 除子与类群}
\label{b3:subsec:2003:02}

每个高度一素理想给出一个整数阶数。把仅有有限多处非零的整数分别指定给这些素理想，就得到一个 Weil 除子；它可以看成有限支撑的阶数资料。同一个非零分式在各处产生的阶数给出主除子。类群所比较的，正是任意这样的整数资料与一个全局分式能够同时产生的资料。

\begin{definition}\label{b3:def:weil-divisor-class-group}
设 \(R\) 为交换 Noether 正规整环，\(K=\operatorname{Frac}(R)\)。对每个高度一素理想 \(\mathfrak p\)，令 \(v_{\mathfrak p}:K^*\to\mathbb Z\) 为离散赋值环 \(R_{\mathfrak p}\) 的归一化离散赋值，即统一化元的值为 \(1\)。以高度一素理想为基的自由交换群
\[
\operatorname{Div}(R)=\bigoplus_{\operatorname{ht}\mathfrak p=1}\mathbb Z[\mathfrak p]
\]
称为 \(R\) 的\term[Weilchuziqun]{Weil 除子群}[Weil divisor group]。其元素 \(D=\sum n_{\mathfrak p}[\mathfrak p]\) 为有限和，称为 Weil 除子；若各 \(n_{\mathfrak p}\ge0\)，则称 \(D\) 为有效除子。对 \(f\in K^*\)，定义主除子
\[
\operatorname{div}(f)=\sum_{\operatorname{ht}\mathfrak p=1}v_{\mathfrak p}(f)[\mathfrak p].
\]
商群
\(\operatorname{Cl}(R)=\operatorname{Div}(R)/\{\operatorname{div}(f):f\in K^*\}\)
称为\term[chuzileiqun]{除子类群}[divisor class group]。两个除子之差为主除子时，称它们线性等价。
\end{definition}
\symidx{\(\operatorname{Div}(R)\)：正规整环的 Weil 除子群}
\symidx{\(\operatorname{div}(f)\)：非零分式 \(f\) 的主除子}
\symidx{\(\operatorname{Cl}(R)\)：正规整环的 Weil 除子类群}

主除子的和确实有限。对 \(0\ne a\in R\)，只须证明：在高度一素理想中，\(v_{\mathfrak p}(a)>0\) 的那些恰好是 \((a)\) 的极小素理想。首先，DVR 中正阶恰表示属于极大理想，故
\(v_{\mathfrak p}(a)>0\) 当且仅当 \(a\in\mathfrak p\)。若 \(\mathfrak q\) 是 \((a)\) 的极小素理想，\cref{b3:thm:principal-ideal-theorem}给出 \(\operatorname{ht}\mathfrak q\le1\)；又因 \(R\) 是整环且 \(a\ne0\)，有 \((0)\subsetneq\mathfrak q\)，故其高度恰为一。反之，若高度一素理想 \(\mathfrak p\) 包含 \((a)\)，可取 \((a)\) 的极小素理想 \(\mathfrak q\subseteq\mathfrak p\)。两者都高度一，若包含严格，则 \((0)\subsetneq\mathfrak q\subsetneq\mathfrak p\) 与 \(\operatorname{ht}\mathfrak p=1\) 矛盾，故 \(\mathfrak p=\mathfrak q\)。Noether 性保证 \((a)\) 只有有限个极小素理想。一般非零分式 \(f=a/b\) 的非零阶数只可能出现在 \((a)\) 或 \((b)\) 的极小素理想处，因此支撑有限。赋值的可加性给出
\(\operatorname{div}(fg)=\operatorname{div}(f)+\operatorname{div}(g)\)，所以主除子确实构成子群。定义及其与 Picard 群的比较可参阅\cite[\S20, pp. 165--167]{Matsumura1989CommutativeRingTheory}。

\ProseParagraphBreak
一个基除子 \([\mathfrak p]\) 在类群中为零，意味着存在 \(f\in K^*\)，使它在 \(\mathfrak p\) 处恰有一阶零，在其他高度一处的阶数都为零。交定理将使 \(f\) 属于 \(R\)，并使 \(\mathfrak p=(f)\)。因此，这个类衡量高度一素理想能否由一个全局方程生成；下面的证明把这种解释与唯一分解判据联系起来。

\begin{theorem}\label{b3:thm:class-zero-ufd}
设 \(R\) 为交换 Noether 正规整环。则 \(R\) 为唯一分解整环，当且仅当 \(\operatorname{Cl}(R)=0\)。
\end{theorem}
\begin{proof}
若 \(R\) 是 UFD，则每个高度一素理想都是 \((\pi)\)，其中 \(\pi\) 为素元，并有 \(\operatorname{div}(\pi)=[(\pi)]\)。因此除子群的每个基元素均为主除子。

反之，若类群为零，给定高度一素理想 \(\mathfrak p\)，存在 \(f\in K^*\) 满足 \(\operatorname{div}(f)=[\mathfrak p]\)。所有赋值非负，所以 \(f\in R\)，且 \(f\in\mathfrak p\)。对任意非零的 \(a\in\mathfrak p\)，在 \(\mathfrak p\) 处有 \(v_{\mathfrak p}(a/f)\ge0\)，在其他高度一素理想处同样非负。交定理给出 \(a/f\in R\)，而零元也属于 \((f)\)，故 \(\mathfrak p=(f)\)。

这样每个高度一素理想都为主理想，\cref{b3:prop:height-one-ufd-criterion}对当前 Noether 整环适用，给出 \(R\) 为 UFD。
\end{proof}

\subsection{Dedekind 理论的高维延伸}
\label{b3:subsec:2003:03}

在 Dedekind 整环中，每个非零素理想都高度一，非零分式理想按赋值唯一分解。这时除子 \(\sum n_{\mathfrak p}[\mathfrak p]\) 对应分式理想 \(\prod\mathfrak p^{n_{\mathfrak p}}\)，本章的类群与\secref{b3:sec:1004}的理想类群一致。

\ProseParagraphBreak
在高维正规整环中，函数的上述恢复性质不能推广到任意理想。例如，设 \(k\) 为域，\(R=k[x,y]\)。理想 \(I=(x,y)\) 在每个高度一素点处均变为整个局部环：高度一素理想不可能同时包含 \(x,y\)。所以 \(I\) 与 \(R\) 有相同的高度一局部数据，但 \(I\ne R\)。特别地，在分式域中有
\[
\bigcap_{\operatorname{ht}\mathfrak p=1}I_{\mathfrak p}
=\bigcap_{\operatorname{ht}\mathfrak p=1}R_{\mathfrak p}
=R\supsetneq I.
\]
函数可以由高度一处的无极点条件恢复，这个普通理想在原点要求的同时消失却无法由那些条件恢复。除子记录的是余维一信息，不记录这个集中在余维二闭点的差别。

\ProseParagraphBreak
要把 Dedekind 分式理想的群运算推广到高维，就应选取恰能由高度一数据恢复的一类模。\secref{b3:sec:2005}将证明，秩一反射模具有这一性质；普通理想乘积则要再取双对偶，才能回到这类模中。

\ProseParagraphBreak
局部化则直接改变哪些高度一素理想仍需记录。逆转一个乘法集后，与它相交的素理想不再对应局部化环的素点；除子群删去相应基元素，类群便按这些类取商。

\begin{proposition}\label{b3:prop:class-group-localization}
设 \(R\) 为交换 Noether 正规整环，\(K=\operatorname{Frac}(R)\)，\(S\subseteq R\setminus\{0\}\) 为乘法集，\(A=S^{-1}R\)。则 \(A\) 仍为 Noether 正规整环，且保留避开 \(S\) 的高度一素理想所诱导的自然群同态
\[
\operatorname{Cl}(R)\longrightarrow\operatorname{Cl}(A)
\]
满射，其核为满足 \(\mathfrak p\cap S\ne\varnothing\) 的高度一素理想 \(\mathfrak p\) 的类所生成的子群。
\end{proposition}
\begin{proof}
局部化保持 Noether 性及整性。\(A\) 的每个素点局部环都同构于 \(R\) 的相应素点局部环，所以正规性按定义保持；两环的分式域都为 \(K\)。素理想对应把 \(\operatorname{Spec}A\) 与 \(R\) 中避开 \(S\) 的素理想对应起来。若 \(\mathfrak p\cap S=\varnothing\)，则 \(\mathfrak p\) 以下的每个素理想也避开 \(S\)，故对应保持止于 \(\mathfrak p\) 的所有素链，给出
\[
\operatorname{ht}_A(S^{-1}\mathfrak p)=\operatorname{ht}_R\mathfrak p.
\]
因此，两边的高度一素理想恰按这个对应匹配，并且
\(A_{S^{-1}\mathfrak p}\cong R_{\mathfrak p}\)，作为 \(K\) 的子环相同。相应归一化赋值也就相同。

在除子群上定义 \(\lambda\)：将避开 \(S\) 的 \([\mathfrak p]\) 送到 \([S^{-1}\mathfrak p]\)，将与 \(S\) 相交的 \([\mathfrak p]\) 送到零。每个 \(A\) 的除子基元素都有这样的原像，所以 \(\lambda\) 满射；其核恰由被删去的基元素生成。赋值相同又使
\[
\lambda\bigl(\operatorname{div}_R(f)\bigr)
=\operatorname{div}_A(f)\qquad(f\in K^*).
\]
这说明 \(\lambda\) 诱导所述类群满射。

若除子 \(D\) 的类落在这个满射的核中，则存在 \(f\in K^*\)，使 \(\lambda(D)=\operatorname{div}_A(f)\)。于是
\(\lambda(D-\operatorname{div}_R(f))=0\)，这个差除子只支撑在与 \(S\) 相交的高度一素理想上。它与 \(D\) 在 \(\operatorname{Cl}(R)\) 中给出同一个类，因此核包含于所述生成子群。反向，这些基除子都被 \(\lambda\) 删除，它们的类均在核中。
\end{proof}

% !TeX root = ../../Book3.tex
% 纲要标题：### 20.4 正规性与正则性的比较
% 纲要位置：计划纲要.md，第 2605 行。

\section{正规性与正则性的比较}
\label{b3:sec:2004}

正规性排除了某些一维缺陷，却不要求每个点正则，更不要求全环唯一分解。二次锥与低维曲线的比较将分别展示这些条件之间成立或失败的蕴含。


% 纲要内容（仅作写作计划，尚非教材正文）：
%
% * 二次锥的高度一正则局部环与闭点嵌入维数；正规不蕴含正则或唯一分解
% * 尖点的整闭性失败与交叉直线的整性失败；二维正规局部环的 CM 性及维数限制
% * 正规化有限性按域上有限型范围说明；消解奇点的几何含义及固有性的后文定义
%

\subsection{正规但非正则的例子}
\label{b3:subsec:2004:01}

\begin{example}\label{b3:exm:normal-quadric-cone}
设 \(k\) 为任意域，
\[
A=k[x,y,z]_{(x,y,z)}/(xy-z^2),\qquad\mathfrak m=(\bar x,\bar y,\bar z).
\]
其中 \(\bar x,\bar y,\bar z\) 分别为 \(x,y,z\) 在商环 \(A\) 中的像。则 \(A\) 是二维正规局部整环，但不是正则局部环。
\end{example}
\begin{proof}
把 \(xy-z^2\) 看成 \(k[y,z][x]\) 中关于 \(x\) 的多项式，其系数 \(y\) 与 \(-z^2\) 互素，所以它是本原多项式。在域 \(k(y,z)\) 上，它是非恒定一次多项式，因而不可约；再由 Gauss 引理，它在 \(k[y,z][x]\) 中也不可约。多项式环 \(k[x,y,z]\) 为 UFD，故不可约元生成素理想，所得商及其局部化 \(A\) 都是整环。\(A\) 由三维正则局部环商去一个非零因子得到，所以维数二且为 CM 环，特别满足 \(S_2\)。

若 \(\mathfrak p\) 高度一，它不能同时包含 \(\bar x,\bar y\)，否则由 \(\bar z^2=\bar x\bar y\) 及素性也包含 \(\bar z\)，便是高度二的 \(\mathfrak m\)。在 \(\bar x\) 可逆处可消去 \(\bar y=\bar z^2/\bar x\)，环成为 \(k[x,x^{-1},z]\) 的局部化；\(\bar y\) 可逆时同样处理，得到 \(k[y,y^{-1},z]\) 的局部化。这些 Laurent 多项式环是 UFD，\cref{b3:prop:height-one-ufd-criterion}使它们的高度一素理想为主理想。局部化对应保持当前素点以下的全部素链，所以它所对应的素理想仍高度一。其局部环的极大理想由一个元素生成，维数也为一，故为正则局部环。高度零处为分式域，所以 \(R_1\) 成立，Serre 判据使 \(A\) 正规。

关系 \(xy-z^2\) 属于 \((x,y,z)^2\)，在原环的 \((x,y,z)/(x,y,z)^2\) 中没有留下任何线性关系。因此 \(\bar x,\bar y,\bar z\) 的一次类在 \(\mathfrak m/\mathfrak m^2\) 中仍线性无关，且生成这个空间，给出
\[
\dim_k\mathfrak m/\mathfrak m^2=3>2=\dim A.
\]
嵌入维数严格大于维数，所以 \(A\) 不正则。
\end{proof}

这个证明不限制特征；即使特征为二，消去计算与余切空间计算都不改变。正规性允许这个余维二处的奇点存在，所以正规局部环未必正则。这个环还不是 UFD：\cref{b3:exm:quadric-nonfactorial-class}将证明高度一素理想 \((\bar x,\bar z)\) 的除子类非零且阶为二。于是同一个二次锥也显示，正规性不能推出唯一分解性。

\subsection{Cohen–Macaulay 但非正规的例子}
\label{b3:subsec:2004:02}

设 \(k\) 为域。一维尖点局部环
\[
R=k[t^2,t^3]_{(t^2,t^3)}\cong k[x,y]_{(x,y)}/(y^2-x^3)
\]
为超曲面环，因而为 CM 环。分式 \(t=y/x\) 满足首一方程 \(T^2-x=0\)，却不在 \(R\) 中：若 \(t=a/b\)，其中 \(a,b\in k[t^2,t^3]\) 且 \(b(0)\ne0\)，则 \(tb\) 的一次项非零，而 \(a\) 的一次项为零，矛盾。因此它不正规。这个计算对任意特征都成立。

\ProseParagraphBreak
交叉直线 \(B=k[x,y]_{(x,y)}/(xy)\) 同样为 CM 超曲面环，但其中 \(\bar x,\bar y\) 都非零，乘积却为零，所以这个局部环不是整环，也就不是正规局部环。尖点破坏的是整环的整闭性，交叉直线在交点处则已经破坏整性。前者的正规化在同一分式域中添加缺失的整元素，后者的正规化把两个分支分开。两者都满足 \(S_2\)，在 Serre 判据中失败的是高度一闭点处的 \(R_1\)。

\ProseParagraphBreak
另一方面，二维交换 Noether 正规局部环 \((R,\mathfrak m)\) 必为 CM 环。Serre 判据给出 \(S_2\)，在闭点处它要求
\[
2=\min\{2,\dim R\}\le\operatorname{depth}R
\le\dim R=2.
\]
深度于是被迫等于维数。这个结论有明确的维数限制：维数至少三时，\(S_2\) 在闭点只要求深度至少二，这还没有达到 CM 性所要求的全部维数。Serre 判据因而不需要假设 CM 性；它所需的高余维深度下界止于二。

\subsection{正规化与奇点消解的区别}
\label{b3:subsec:2004:03}

正规化将整环补成其分式域中的整闭环；对约化 Noether 环，则分别取各极小素理想对应的整环商的整闭包，再取它们的直积，得到各不可约分量的正规化。它针对首一方程，而正则性针对局部极大理想的生成结构。对域上有限型代数，\cref{b3:thm:normalization-finite}已证明这项整扩张有限；一般 Noether 环的正规化不能仅由定义保证有限。\cref{b3:exm:normal-quadric-cone}的环已经正规，其正规化就是自身，原点的非正则性却仍然存在。

\ProseParagraphBreak
这里先说明奇点消解的几何含义，概形的构造见\secref{b3:sec:2401}，固有态射的正式定义见\cref{b3:def:separated-finite-type-proper}。对于域上整的有限型代数簇，\term[qidianxiaojie]{奇点消解}[resolution of singularities]寻求一个固有双有理态射 \(Y\to X\)，使源空间 \(Y\) 的每个局部环都正则。双有理性在这里指该态射在源与目标的适当稠密开子集上给出同构，因而诱导相同的有理函数域；固有性还对态射的有限型、分离性及基变换后的闭性提出要求。这一过程允许改变几何空间，以替换原来的奇点。这里的“消解”是一种几何操作，与本编用自由模构造自由消解的含义不同。

\ProseParagraphBreak
曲线容易掩盖这个区别：一维 Noether 正规局部整环是 DVR，所以曲线的正规化在局部正则性上已经起了很大作用；二维锥面则直接显示，正规之后仍可能需要改变几何空间。还应区分正则与基域上的光滑，非完美基域上的差别已在\secref{b3:sec:1805}讨论。

% !TeX root = ../../Book3.tex
% 纲要标题：### 20.5* 秩一反射模、除子类群与局部阶乘性
% 纲要位置：计划纲要.md，第 2611 行。

\OptionalSection{秩一反射模、除子类群与局部阶乘性}{b3:sec:2005}

对无挠模取两次对偶，会保留所有高度一素点能够检测到的元素。秩一时，这种反射化与除子的阶数描述相合，从而把模同构类、除子类群与局部阶乘性联系起来。


% 纲要内容（仅作写作计划，尚非教材正文）：
%
% * 在 Noether 正规整环上定义双对偶及反射性；先直接计算 (x,y) 的双对偶，再证明共同分式向量空间内的高度一交公式、无挠模反射性与 S_2 的等价以及分式理想的理想商表达
% * 证明双对偶张量积及对偶的有限生成、秩一与反射性；以反射分式理想解释除子类群运算
% * 保留理想正号约定，并说明与几何中 O(D) 约定的负号差；比较可逆模的 Picard 群与 Weil 除子类群
% * 用第18.4节的唯一分解性解释正则局部情形的类群消失；计算二次锥的阶二非零类及普通平方与反射平方的区别
% * 区分局部阶乘、全局主理想性和唯一分解；保持仿射环论范围，完整层论表述留给第24章之后的几何教材
%
% ---
%

\subsection{双对偶、反射性与高度一局部数据}
\label{b3:subsec:2005:01}

\begin{definition}\label{b3:def:reflexive-module}
设 \(R\) 为交换 Noether 正规整环，\(M\) 为有限生成的 \(R\) 模。令
\(M^*=\operatorname{Hom}_R(M,R)\)、\(M^{**}=\operatorname{Hom}_R(M^*,R)\)。若求值映射
\[
M\longrightarrow M^{**},\qquad m\longmapsto\bigl(\varphi\mapsto\varphi(m)\bigr)
\]
为同构，则称 \(M\) 为\term[fanshemo]{反射模}[reflexive module]。若再有
\(\dim_K(M\otimes_RK)=1\)，其中 \(K=\operatorname{Frac}(R)\)，则称其秩为一。
\end{definition}
\symidx{\(M^*\)、\(M^{**}\)：模的对偶与双对偶}

有限生成投射模是反射模，因为有限秩自由模的求值映射为同构，这一性质可取直和项。有限生成无挠模的求值映射则至少为单射。事实上，将 \(M\) 嵌入 \(V=M\otimes_RK\) 后，对任意 \(0\ne m\in M\)，选取 \(K\) 线性函数 \(\lambda:V\to K\) 使 \(\lambda(m)\ne0\)。设 \(m_1,\ldots,m_s\) 生成 \(M\)，清掉有限个 \(\lambda(m_i)\) 的分母，得到 \(0\ne a\in R\) 使 \(\varphi=a\lambda\) 在 \(M\) 上取值于 \(R\)。于是 \(\varphi\in M^*\)，且 \(\varphi(m)\ne0\)，所以 \(m\) 的求值像不为零。有限生成性正用于选取这个共同分母。

\ProseParagraphBreak
单射未必满射。设 \(k\) 为域，\(R=k[x,y]\)、\(I=(x,y)\)。对任意 \(\varphi:I\to R\)，记 \(a=\varphi(x)\)、\(b=\varphi(y)\)。生成元的关系给出 \(ya=xb\)。模掉 \(x\) 后，\(k[y]\) 中的 \(y\) 为非零因子，故 \(a=xc\) 对某个 \(c\in R\) 成立；在整环中约去 \(x\)，再得 \(b=yc\)。因此每个 \(\varphi\) 都是乘以唯一的 \(c\)，即 \(I^*\cong R\)，以包含映射 \(I\hookrightarrow R\) 为基。再次对偶得到 \(I^{**}\cong R\)，而求值映射把 \(f\in I\) 送到 \(f\in R\)，故恰为严格包含 \(I\subsetneq R\)。商 \(R/I\cong k\) 只支撑在原点 \((x,y)\)。这一处的缺失不能由高度一局部化检测：高度一素理想不可能同时包含 \(x,y\)，所以每个 \(I_{\mathfrak p}\) 都等于 \(R_{\mathfrak p}\)。

\ProseParagraphBreak
对一般有限生成无挠模，局部化的映射都为单射，可以把
\[
M\subseteq M_{\mathfrak p}\subseteq V=M\otimes_RK
\]
放在同一个 \(K\) 向量空间中。下面的交集就是 \(V\) 内这些子模的普通交集。证明将把“所有线性函数取值在 \(R\) 中”逐点化，再用高度一 DVR 上无挠模自由的性质识别双对偶。

\begin{theorem}\label{b3:thm:reflexive-height-one}\label{b3:thm:rank-one-dual-colon}\label{b3:thm:reflexive-serre-s2}
设 \(R\) 为交换 Noether 正规整环，\(K=\operatorname{Frac}(R)\)，\(M\) 为有限生成的无挠 \(R\) 模，\(V=M\otimes_RK\)，\(M^{**}=\operatorname{Hom}_R(\operatorname{Hom}_R(M,R),R)\)。在 \(V\) 中自然识别后，有
\[
M^{**}=\bigcap_{\operatorname{ht}\mathfrak p=1}M_{\mathfrak p}.
\]
这里交集遍历 \(R\) 的所有高度一素理想 \(\mathfrak p\)；若 \(R\) 为域，空交集按 \(V\) 解释。
特别地，对这个无挠模 \(M\)，反射性、等于右侧交集及满足 Serre 条件 \((S_2)\) 三个条件等价；\(M^{**}\) 总是反射模。这里 \((S_2)\) 指对每个 \(\mathfrak p\in\operatorname{Supp}_R M\)，有
\[
\operatorname{depth}_{R_{\mathfrak p}}M_{\mathfrak p}
\ge\min\{2,\dim_{R_{\mathfrak p}}M_{\mathfrak p}\}.
\]

对非零分式理想 \(I\subset K\)，把到 \(R\) 的同态识别为 \(K\) 中的标量乘法后，还有
\[
I^*=(R:I)=\{a\in K:aI\subseteq R\},\qquad
I^{**}=(R:(R:I)).
\]
若 \(k\) 为域、\(R=k[x,y]\)、\(I=(x,y)^2\)，则 \(I^{**}=R\)，求值映射为严格包含 \(I\subsetneq R\)。
\end{theorem}
\begin{proof}
取 \(M\) 的一组有限生成元，按各生成元上的取值把 \(M^*\) 嵌入某个有限秩自由模。Noether 性使 \(M^*\) 有限生成；再次对偶，同样得到 \(M^{**}\) 有限生成。它们也都是无挠模，且有限生成模在 Noether 环上有限展示，所以\cref{b3:thm:hom-localization}给出 \(M^*\otimes_RK\cong V^*\)，再对偶得到 \(M^{**}\otimes_RK\cong V\)。求值意义下，\(v\in V\) 属于 \(M^{**}\)，当且仅当对所有 \(\varphi\in M^*\)，都有 \(\varphi(v)\in R\)。利用\cref{b3:thm:normal-height-one-intersection}，这等价于对每个高度一 \(\mathfrak p\)，所有这些值均属于 \(R_{\mathfrak p}\)。

又由\cref{b3:thm:hom-localization}，\((M^*)_{\mathfrak p}\cong(M_{\mathfrak p})^*\)，所以这个条件等价于 \(v\in(M_{\mathfrak p})^{**}\)。\(R_{\mathfrak p}\) 是 DVR，\cref{b3:prop:dvr-order-ideals}说明它为主理想整环；第二卷定理4.3.2“有限生成模的不变因子分解”使有限生成无挠模 \(M_{\mathfrak p}\) 自由。因此 \((M_{\mathfrak p})^{**}=M_{\mathfrak p}\)，得到交集公式。

最后，\(M^{**}\) 有限生成且无挠，且其高度一局部化仍为 \(M_{\mathfrak p}\)。把已证交集公式用于 \(M^{**}\)，即得 \((M^{**})^{**}=M^{**}\)。

为联系 \((S_2)\)，零模的情形显然，以下设 \(M\ne0\)。先证明 \(M^{**}\) 满足这一深度条件。对有限生成模 \(M^*\) 取有限展示 \(R^b\to R^a\to M^*\to0\)，对偶后得到
\[
0\longrightarrow M^{**}\longrightarrow R^a\longrightarrow R^b.
\]
记最后一条箭头的像为 \(C\)，则 \(C\) 是有限生成无挠模。\(M^{**}\) 也无挠，所以在高度一素点处，任一非零的局部极大理想元素都在它上面单射，深度至少为一。在高度至少二的 \(\mathfrak p\) 处，\cref{b3:thm:serre-normality}使 \(\operatorname{depth}R_{\mathfrak p}\ge2\)，而 \(C_{\mathfrak p}\) 无挠使其深度至少为一；若该模为零，采用深度为无穷大的约定。对局部化后的短正合列 \(0\to M^{**}\to R^a\to C\to0\) 应用\cref{b3:thm:depth-lemma}，得到 \(\operatorname{depth}M^{**}_{\mathfrak p}\ge2\)。非零无挠模的支集是整个素谱，因而局部模维数等于 \(\dim R_{\mathfrak p}\)；高度零处无额外要求。于是 \(M^{**}\) 满足 \((S_2)\)，反射的 \(M\) 也满足它。

反过来，设无挠模 \(M\) 满足 \((S_2)\)，并令 \(Q=M^{**}/M\)。交集公式及分式域上的识别说明，\(Q\) 在高度零和一的点都为零。若 \(Q\ne0\)，在其支集中取极小素理想 \(\mathfrak p\)，则 \(\operatorname{ht}\mathfrak p\ge2\)，且 \(Q_{\mathfrak p}\) 是非零有限长度模，深度为零。但局部化后的短正合列给出
\[
\operatorname{depth}_{R_{\mathfrak p}}Q_{\mathfrak p}
\ge\min\{\operatorname{depth}_{R_{\mathfrak p}}M_{\mathfrak p}-1,
          \operatorname{depth}_{R_{\mathfrak p}}M^{**}_{\mathfrak p}\}
\ge1,
\]
矛盾。因此 \(Q=0\)，求值映射为同构。

对分式理想 \(I\)，有 \(I\otimes_RK\cong K\)。任意 \(R\) 线性同态 \(I\to R\) 张量到 \(K\) 后，成为 \(K\to K\) 的线性函数，故为乘以某个 \(a\in K\)。它在 \(I\) 上取值于 \(R\)，恰好要求 \(aI\subseteq R\)；反之，每个满足此条件的标量都给出同态。因此 \(I^*=(R:I)\)。再对这个秩一无挠模应用同一识别，就得到 \(I^{**}=(R:(R:I))\)，且求值映射对应 \(I\subseteq I^{**}\) 的自然包含。

在 \(R=k[x,y]\) 中，高度一素理想不能同时包含 \(x,y\)，故 \(((x,y)^2)_{\mathfrak p}=R_{\mathfrak p}\) 对每个高度一 \(\mathfrak p\) 成立。交集公式遂给出 \(((x,y)^2)^{**}=R\)。又因 \(1\notin(x,y)^2\)，求值映射不满。\(R/(x,y)^2\) 只支撑在高度二的原点，并且保留常数项和一次项；双对偶同样把这种只在高余维处的缺失补成整个环。
\end{proof}

任意有限生成模的挠子模会被所有到 \(R\) 的同态零化，因此双对偶实际等于其无挠商的双对偶。又因为双对偶无挠，反射模必然无挠。所以上一定理对 Noether 正规整环上的任意有限生成模给出：反射性等价于无挠且满足 \((S_2)\)。无挠使原模能嵌入共同的分式向量空间，\((S_2)\) 则排除高度一局部化无法检测的商 \(M^{**}/M\)。例如前面的 \((x,y)\subset k[x,y]\) 在原点的深度只有一，正是它不反射的原因。

\subsection{分式理想的反射包络、Picard 群与除子类群}
\label{b3:subsec:2005:02}

秩一有限生成无挠模选定 \(M\otimes_RK\cong K\) 后成为一个非零分式理想 \(I\subset K\)。这里“分式理想”指存在 \(0\ne a\in R\) 使 \(aI\subseteq R\) 的非零 \(R\) 子模；Noether 性保证它有限生成。对每个高度一 \(\mathfrak p\)，定义
\[
v_{\mathfrak p}(I)=\min\{v_{\mathfrak p}(f):0\ne f\in I\}.
\]
取有限组生成元即可求此最小值，且只有有限多个素理想处的值非零。

\begin{lemma}\label{b3:lem:divisor-fractional-module}
设 \(R\) 为交换 Noether 正规整环，\(K\) 为其分式域。对每个高度一素理想 \(\mathfrak p\)，记 \(v_{\mathfrak p}:K^*\rightarrow\mathbb Z\) 为 \(R_{\mathfrak p}\) 的归一化离散赋值，即统一化元的值为 \(1\)。给定整数系数仅有有限个非零的除子
\(D=\sum n_{\mathfrak p}[\mathfrak p]\)，令
\[
I_D=\{0\}\cup\{f\in K^*:v_{\mathfrak p}(f)\ge n_{\mathfrak p}
\;\text{对所有高度一}\;\mathfrak p\}.
\]
则 \(I_D\) 是非零有限生成分式理想，且
\((I_D)_{\mathfrak p}=\pi_{\mathfrak p}^{n_{\mathfrak p}}R_{\mathfrak p}\)，其中 \(\pi_{\mathfrak p}\) 为任一局部统一化元。因此 \(I_D\) 反射。
\end{lemma}
\begin{proof}
给定各高度一素点的阶数下界后，先要找出一个满足全部下界的非零元素，并用共同分母把整个模放回 \(R\)。随后还须在每个指定素点实现精确阶数，才能证明局部化等式，最后由交集公式得到反射性。

先验证 \(I_D\) 是 \(R\) 子模。对非零的 \(f,g\in I_D\)，若 \(f+g\ne0\)，赋值不等式给出
\[
v_{\mathfrak p}(f+g)\ge
\min\{v_{\mathfrak p}(f),v_{\mathfrak p}(g)\}\ge n_{\mathfrak p}.
\]
每个 \(r\in R\setminus\{0\}\) 都满足 \(v_{\mathfrak p}(r)\ge0\)，所以 \(rf\in I_D\)；含零的情形也成立。对有限多个满足 \(n_{\mathfrak p}>0\) 的素理想，分别选 \(0\ne b_{\mathfrak p}\in\mathfrak p\)。由于 \(v_{\mathfrak p}(b_{\mathfrak p})\ge1\)，乘积
\[
h=\prod_{n_{\mathfrak p}>0}b_{\mathfrak p}^{n_{\mathfrak p}}
\]
在每个正系数素点处都有足够大的阶数，在其余素点处阶数非负，所以 \(0\ne h\in I_D\)。一个因子可能同时提高别的素点的阶数，但所有要求都是下界，额外提高不会破坏任何条件；若无正系数，取空乘积 \(h=1\)。

对有限多个满足 \(n_{\mathfrak p}<0\) 的素理想作同样选择，令
\[
a=\prod_{n_{\mathfrak p}<0}b_{\mathfrak p}^{-n_{\mathfrak p}}\ne0.
\]
则对每个 \(f\in I_D\setminus\{0\}\)，\(af\) 的所有高度一赋值均非负。由\cref{b3:thm:normal-height-one-intersection}，\(aI_D\subseteq R\)。它是 \(R\) 的非零理想，Noether 性使其有限生成；乘以 \(a\) 为模同构，所以 \(I_D\) 也是有限生成分式理想。

固定高度一素理想 \(\mathfrak p\)。先取 \(f\in K^*\) 满足 \(v_{\mathfrak p}(f)=n_{\mathfrak p}\)，例如取局部统一化元的相应整数次幂。\(D\) 与 \(\operatorname{div}(f)\) 均有有限支撑，故只有有限多个其他素理想 \(\mathfrak q\) 不满足 \(v_{\mathfrak q}(f)\ge n_{\mathfrak q}\)。对每个这样的 \(\mathfrak q\)，选 \(b_{\mathfrak q}\in\mathfrak q\setminus\mathfrak p\)；不同的高度一素理想不能相互包含，所以这样的元素存在。此时
\[
v_{\mathfrak q}(b_{\mathfrak q})>0,
\qquad v_{\mathfrak p}(b_{\mathfrak q})=0.
\]
取足够大的整数 \(e_{\mathfrak q}\ge0\)，使 \(g=f\prod b_{\mathfrak q}^{e_{\mathfrak q}}\) 修补所有不足的阶数。每个修补因子在 \(\mathfrak p\) 处为单位，因而仍有 \(v_{\mathfrak p}(g)=n_{\mathfrak p}\)；在其他素点处它的阶数非负，故不会制造新的不足。于是 \(g\in I_D\)。所有 \(I_D\) 中元素在 \(\mathfrak p\) 处的阶数均不小于 \(n_{\mathfrak p}\)，而 \(g\) 恰好达到该下界，所以
\[
(I_D)_{\mathfrak p}=gR_{\mathfrak p}
=\pi_{\mathfrak p}^{n_{\mathfrak p}}R_{\mathfrak p}.
\]
这些局部化在 \(K\) 中的交集，按赋值不等式正是 \(I_D\)。因此\cref{b3:thm:reflexive-height-one}给出 \(I_D^{**}=I_D\)。
\end{proof}

相加除子要求各局部阶数相加，对模的对应运算因而从张量积开始。不过，两个秩一反射模的普通张量积可能带有挠，也可能在高余维处不反射。先去挠再取双对偶，才得到由高度一阶数唯一确定的模。后面的\cref{b3:exm:quadric-nonfactorial-class}将给出具体计算：普通理想平方 \(\mathfrak p^2=\bar x\mathfrak m\) 严格小于 \((\bar x)\)，而反射化后两者相同。

\begin{theorem}\label{b3:thm:class-group-reflexive}
设 \(R\) 为交换 Noether 正规整环，\(K=\operatorname{Frac}(R)\)。对 \(R\) 模 \(X\) 记 \(X^*=\operatorname{Hom}_R(X,R)\)、\(X^{**}=\operatorname{Hom}_R(X^*,R)\)。对高度一素理想 \(\mathfrak p\) 记 \(v_{\mathfrak p}\) 为 \(R_{\mathfrak p}\) 的归一化离散赋值；对非零分式理想 \(I\subset K\)，记 \(v_{\mathfrak p}(I)=\min\{v_{\mathfrak p}(f):0\ne f\in I\}\)。秩一有限生成反射 \(R\) 模的同构类在
\[
[M]\mathbin{\star}[N]=[(M\otimes_RN)^{**}]
\]
下组成交换群，其中 \(M,N\) 为任意两个秩一有限生成反射 \(R\) 模，单位为 \([R]\)，\([M]\) 的逆为 \([M^*]\)。对这样的模选取非零分式理想表示 \(I\subset K\) 后，映射
\[
[I]\longmapsto\left[\sum_{\operatorname{ht}\mathfrak p=1}
v_{\mathfrak p}(I)[\mathfrak p]\right]
\]
给出该群到 \(\operatorname{Cl}(R)\) 的同构。
\end{theorem}
\begin{proof}
要使用高度一交公式，先须把张量积中的挠去掉。令 \(T\) 为 \(M\otimes_RN\) 的挠子模，则无挠商 \(L=(M\otimes_RN)/T\) 有限生成，且
\[
L\otimes_RK\cong
(M\otimes_RK)\otimes_K(N\otimes_RK)\cong K.
\]
若 \(t\in T\)，取 \(0\ne a\in R\) 使 \(at=0\)，则对任意 \(\varphi:M\otimes_RN\to R\)，有 \(a\varphi(t)=0\)。\(R\) 为整环，故 \(\varphi(t)=0\)，所以每个这样的同态都唯一经过 \(L\) 分解。因而 \((M\otimes_RN)^{**}\cong L^{**}\)。\Cref{b3:thm:reflexive-height-one}现在适用于 \(L\)，保证 \(L^{**}\) 有限生成、反射且仍为秩一，这就验证了乘积的闭性。

对偶将给出逆元，因此还要验证 \(M^*\) 属于同一类模。取 \(V=M\otimes_RK\)，将 \(M^*\) 视为 \(V^*\) 的子模。由有限组生成元的取值，\(M^*\) 嵌入某个有限秩自由模，Noether 性使它有限生成；局部化到 \(K\) 后有 \(M^*\otimes_RK\cong V^*\)，所以秩为一。对 \(\varphi\in V^*\)，\cref{b3:thm:hom-localization}与高度一交定理给出
\[
\begin{aligned}
\varphi\in\bigcap_{\operatorname{ht}\mathfrak p=1}(M^*)_{\mathfrak p}
&\Longleftrightarrow
\varphi(M)\subseteq R_{\mathfrak p}
\quad\text{对每个高度一}\;\mathfrak p\\
&\Longleftrightarrow \varphi(M)\subseteq R
\Longleftrightarrow \varphi\in M^*.
\end{aligned}
\]
因此 \(M^*\) 等于自身的高度一局部化之交，由\cref{b3:thm:reflexive-height-one}可知它反射。

群运算可以通过局部阶数核验。选定分式理想表示 \(I,J\subset K\) 后，乘法映射 \(I\otimes_RJ\rightarrow IJ\) 为满射，且张量到 \(K\) 后为同构。因此其核的每个元素都被某个非零标量零化；反之，目标 \(IJ\subset K\) 无挠，所有挠元素都落在核中，故该核恰为挠子模。于是双对偶乘积可写成 \((IJ)^{**}\)。\Cref{b3:thm:rank-one-dual-colon}也把 \(I^*\) 识别为 \((R:I)\)。在高度一离散赋值环 \(R_{\mathfrak p}\) 中，写
\[
I_{\mathfrak p}=\pi_{\mathfrak p}^{a}R_{\mathfrak p},
\qquad J_{\mathfrak p}=\pi_{\mathfrak p}^{b}R_{\mathfrak p}.
\]
普通乘积局部化为 \(\pi_{\mathfrak p}^{a+b}R_{\mathfrak p}\)，其双对偶不改变这个自由模；对偶局部化则为 \(\pi_{\mathfrak p}^{-a}R_{\mathfrak p}\)。所以
\[
v_{\mathfrak p}((IJ)^{**})
=v_{\mathfrak p}(I)+v_{\mathfrak p}(J),
\qquad
v_{\mathfrak p}(I^*)=-v_{\mathfrak p}(I).
\]
反射分式理想由这些局部化唯一确定。三个理想的两种结合次序在每个高度一素点处都具有三个阶数之和，因而给出同一个反射分式理想。交换律来自阶数相加的交换性，\(R\) 的全部阶数为零使它成为单位，而 \(I\) 与 \(I^*\) 的阶数之和均为零，故 \((II^*)^{**}=R\)。

现在验证到除子类群的映射。换一个 \(I\otimes_RK\cong K\) 的识别，相当于把 \(I\) 乘以 \(f\in K^*\)，其除子增加 \(\operatorname{div}(f)\)，所以除子类不变。更一般地，分式理想之间的模同构张量到 \(K\) 后必为乘以一个非零标量，因此同构的模得到相同的除子类。这证明映射在模同构类上良定义；已证的阶数加法公式也说明它是群同态。

若 \(I,J\) 的除子类相同，则存在 \(f\in K^*\)，使
\[
v_{\mathfrak p}(J)=v_{\mathfrak p}(I)+v_{\mathfrak p}(f)
=v_{\mathfrak p}(fI)
\]
对每个高度一 \(\mathfrak p\) 成立。DVR 中的分式理想由阶数唯一确定，故 \(J_{\mathfrak p}=(fI)_{\mathfrak p}\)。两边都是反射分式理想，交集公式给出 \(J=fI\)，从而 \(I\cong J\)，证明单射。

最后，对任意除子 \(D\)，\cref{b3:lem:divisor-fractional-module}给出的 \(I_D\) 是秩一反射分式理想，且在每个高度一素点处具有 \(D\) 所指定的精确阶数。因此其模类映到 \([D]\)，证明满射。
\end{proof}

这里采用理想类群的正号约定。对高度一素理想 \(\mathfrak p\)，\(I_{[\mathfrak p]}\) 中元素在所有高度一素点的阶数都非负，故它们属于 \(R\)；再要求 \(v_{\mathfrak p}(f)\ge1\)，就等价于 \(f\in\mathfrak pR_{\mathfrak p}\cap R=\mathfrak p\)。因此 \(I_{[\mathfrak p]}=\mathfrak p\)，它是反射理想，其模类送到除子类 \([\mathfrak p]\)。几何文献通常以 \(\operatorname{div}(f)+D\ge0\) 为 \(\mathcal O(D)\) 的非零有理函数截面条件；若 \(D=\sum n_{\mathfrak p}[\mathfrak p]\)，这要求 \(v_{\mathfrak p}(f)\ge-n_{\mathfrak p}\)。因此在仿射情形，\(I_D\) 对应 \(\mathcal O(-D)\) 的全局截面；本节的模类到除子类映射，与几何中按 \(L\simeq\mathcal O(D)\) 取除子类的约定相差一个负号\cite[Chapter II, \S6, pp. 144--145]{Hartshorne1977}。

\begin{proposition}\label{b3:prop:picard-class-group}
设 \(R\) 为交换 Noether 正规整环。按照\cref{b3:thm:class-group-reflexive}的理想类群符号约定，把秩一有限生成投射模视为秩一反射模，得到规范单射
\(\operatorname{Pic}(R)\hookrightarrow\operatorname{Cl}(R)\)，其中 \(\operatorname{Pic}(R)\) 是秩一有限生成投射 \(R\) 模的同构类在张量积下所成的群。其像恰由那些在每个素点处都是秩一自由模的反射模类组成。
\end{proposition}
\begin{proof}
有限生成投射模反射，两个这样的模的张量积仍投射，故上述群运算与 Picard 群运算一致。如果其类在 \(\operatorname{Cl}(R)\) 中为零，上一定理使该模同构于 \(R\)，证明单射。由于 \(R\) 为 Noether 环，此处有限生成模都是有限展示模，\cref{b3:thm:flatness-local}、\cref{b3:thm:finite-presentation-flat-projective}和\cref{b3:lem:finite-flat-local-free}使逐点自由秩一等价于秩一有限生成投射性。
\end{proof}

这个单射比较的是两种局部要求：秩一反射模能由全部高度一局部化恢复，但不必在高余维素点处自由；秩一投射模则要求在所有素点处都自由秩一。前者给出除子类群，后者才给出 Picard 群。

\subsection{正则局部环的类群消失与非阶乘模型}
\label{b3:subsec:2005:03}

\cref{b3:thm:regular-local-ufd}与\cref{b3:thm:class-zero-ufd}说明，正则局部环的除子类群为零。这个结论用到了比正规性更强的假设；二维锥面能给出一个非零类。

\begin{example}\label{b3:exm:quadric-nonfactorial-class}
设 \(k\) 为域，取\cref{b3:exm:normal-quadric-cone}中的环
\(A=k[x,y,z]_{(x,y,z)}/(xy-z^2)\)，令 \(\mathfrak m=(\bar x,\bar y,\bar z)\)、\(\mathfrak p=(\bar x,\bar z)\)。则 \(\mathfrak p\) 高度一，类 \([\mathfrak p]\in\operatorname{Cl}(A)\) 非零，且其阶为二。
其最少生成元个数为 \(\mu_A(\mathfrak p)=2\)，而普通平方与反射平方分别满足
\[
\mathfrak p^2=\bar x\mathfrak m\subsetneq\bar xA,
\qquad(\mathfrak p^2)^{**}=\bar xA.
\]
\end{example}
\begin{proof}
商 \(A/\mathfrak p\cong k[y]_{(y)}\) 为整环，所以 \(\mathfrak p\) 为素理想。又有
\[
A/(\bar x)\cong k[y,z]_{(y,z)}/(z^2),
\]
右边的唯一极小素理想为 \((\bar z)\)，对应到 \(A\) 正是 \(\mathfrak p\)。主理想定理遂给出 \(\operatorname{ht}\mathfrak p=1\)。在 \(A_{\mathfrak p}\) 中，\(\bar y\) 为单位，\(\bar x=\bar z^2/\bar y\)，故 \(\bar z\) 生成极大理想，\(v_{\mathfrak p}(\bar x)=2\)。其他高度一素理想不含 \(\bar x\)，所以
\[
\operatorname{div}(\bar x)=2[\mathfrak p],
\qquad 2[\mathfrak p]=0\quad\text{在}\;\operatorname{Cl}(A)\;\text{中}.
\]
这先说明该类的阶整除二，还须排除它为零。

由前面 \(I_{[\mathfrak p]}=\mathfrak p\) 的识别，\(\mathfrak p\) 为秩一反射模。若 \([\mathfrak p]=0\)，\cref{b3:thm:class-group-reflexive}使 \(\mathfrak p\cong A\)，于是它应由 \(1\) 的像主生成。但 \(\bar x,\bar z\) 在 \(\mathfrak p/\mathfrak m\mathfrak p\) 中线性独立：若非平凡常系数线性组合落在 \(\mathfrak m\mathfrak p\subseteq\mathfrak m^2\)，就与\cref{b3:exm:normal-quadric-cone}中 \(\mathfrak m/\mathfrak m^2\) 的三个坐标线性独立矛盾。再由\cref{b3:cor:local-generators}，
\[
\mu_A(\mathfrak p)=\dim_k\mathfrak p/\mathfrak m\mathfrak p=2.
\]
因此 \(\mathfrak p\) 非主，\([\mathfrak p]\ne0\)。非零与两倍为零合在一起，才说明其阶恰为二。

这个理想的普通平方还保留一个只在闭点处的缺失。由 \(\bar z^2=\bar x\bar y\)，有
\[
\mathfrak p^2=(\bar x^2,\bar x\bar z,\bar z^2)
=\bar x\mathfrak m\subsetneq(\bar x),
\qquad (\mathfrak p^2)^{**}=(\bar x).
\]
严格包含是因为若 \(\bar x\in\bar x\mathfrak m\)，整环中约去非零的 \(\bar x\) 就会得到 \(1\in\mathfrak m\)。另一方面，在每个高度一素理想 \(\mathfrak q\) 处，都有 \(\mathfrak mA_{\mathfrak q}=A_{\mathfrak q}\)：\(\mathfrak q\ne\mathfrak m\)，故 \(\mathfrak m\) 中有元素在那里成为单位。于是 \((\bar x\mathfrak m)_{\mathfrak q}=\bar xA_{\mathfrak q}\)，高度一交公式将其双对偶补成主理想 \((\bar x)\)。
\end{proof}

这里 \(\operatorname{Pic}(A)=0\)，因为局部环上的有限生成投射模自由；但 \(\operatorname{Cl}(A)\ne0\)。所以 Picard 群到类群的单射一般不是满射：\(\mathfrak p\) 能由高度一数据恢复，却在闭点需要两个生成元，不能在该点自由秩一。

\ProseParagraphBreak
乘法满射 \(\mathfrak p\otimes_A\mathfrak p\to\mathfrak p^2\) 在分式域上是同构，核恰为挠子模；去挠后得到 \(\mathfrak p^2\)，再取双对偶补成 \((\bar x)\)。因此\cref{b3:thm:class-group-reflexive}中的 \(2[\mathfrak p]=0\) 对应反射化后的平方为主，并不表示普通理想平方已是主理想。普通平方在闭点缺少的部分由
\(\bar xA/\bar x\mathfrak m\cong A/\mathfrak m\) 给出，正是双对偶补回的余维二差别。

\subsection{局部阶乘性、全局主理想性与唯一分解}
\label{b3:subsec:2005:04}

\begin{definition}\label{b3:def:locally-factorial-domain}
设 \(R\) 为交换 Noether 整环。若每个素理想 \(\mathfrak q\) 处的局部环 \(R_{\mathfrak q}\) 都为唯一分解整环，则称 \(R\) 为\term[jubujiechengzhenghuan]{局部阶乘整环}[locally factorial domain]。
\end{definition}

局部阶乘性要使每个高度一素理想在任意素点附近都能由一个元素生成。用前面的模描述，这等价于每个除子类都能由一个逐点自由秩一的模表示，正是 Picard 群到除子类群的单射满射。

\begin{proposition}\label{b3:prop:locally-factorial-pic-cl}
设 \(R\) 为交换 Noether 正规整环。则 \(R\) 局部阶乘，当且仅当
本节的规范单射 \(\operatorname{Pic}(R)\hookrightarrow\operatorname{Cl}(R)\) 为同构。
\end{proposition}
\begin{proof}
设 \(R\) 局部阶乘。每个高度一素理想 \(\mathfrak p\) 的局部化，在不含它的素点处为单位理想，在包含它的素点处是 UFD 的高度一素理想，因而主生成。它是有限生成模，且在每个素点处自由秩一，所以为可逆理想。其除子恰为 \([\mathfrak p]\)，这些元素生成类群，故 Picard 群映射满射，结合已证单射即为同构。

反之，假定映射满射。对高度一素理想 \(\mathfrak p\)，有 \(\mathfrak pR_{\mathfrak p}\cap R=\mathfrak p\)，且其余高度一素点处的局部化是单位理想。结合高度一交定理，可得
\[
\mathfrak p=\{0\}\cup\{f\in R\setminus\{0\}:v_{\mathfrak p}(f)\ge1\}
=I_{[\mathfrak p]}.
\]
所以 \(\mathfrak p\) 是反射理想。它的类有一个可逆代表，\cref{b3:thm:class-group-reflexive}使 \(\mathfrak p\) 与该代表同构，故 \(\mathfrak p\) 也是可逆模。固定 \(\mathfrak q\)，\(R_{\mathfrak q}\) 的每个素理想都唯一来自某个 \(\mathfrak p\subseteq\mathfrak q\)。位于 \(\mathfrak p\) 以下的每个素理想也包含于 \(\mathfrak q\)，因而都避开乘法集 \(R\setminus\mathfrak q\)；局部化的素理想对应完整保留了以 \(\mathfrak p\) 为终点的全部素链。于是
\[
\operatorname{ht}_{R_{\mathfrak q}}(\mathfrak pR_{\mathfrak q})
=\operatorname{ht}_R\mathfrak p.
\]
特别地，高度一素理想恰好来自高度一的 \(\mathfrak p\subseteq\mathfrak q\)。由于 \(\mathfrak p\) 逐点自由秩一，其局部化 \(\mathfrak pR_{\mathfrak q}\) 主生成。Noether 整环中由高度一素理想主性推出 UFD 的论证已在\cref{b3:thm:class-zero-ufd}证明，遂知每个 \(R_{\mathfrak q}\) 均为 UFD。
\end{proof}

局部阶乘性允许不同开邻域采用不同生成元；全局 UFD 则要求每个高度一素理想在整个环中主生成；主理想整环更要求每个理想都主生成。Dedekind 整环总是局部阶乘，但可以有非零理想类群，例如\secref{b3:sec:1004}中的 \(\mathbb Z[\sqrt{-5}]\)。这个环因而不是 UFD。反过来，对任意域 \(k\)，\(k[x,y]\) 为 UFD，其中的 \((x,y)\) 却不是主理想，所以它不是主理想整环。逐点主生成高度一素理想、全局主生成高度一素理想、全局主生成所有理想，是三个不同的要求。

\ProseParagraphBreak
下面的图先保留具体除子和反射分式理想在 \(K\) 中的位置，随后再取等价类。第一行的对应是 \(D\mapsto I_D\)，逆向从具体理想取各高度一阶数；将理想乘以 \(f\in K^*\) 时，其除子增加 \(\operatorname{div}(f)\)。取主除子和非零标量倍乘的商以后，才得到与嵌入无关的模同构类。

\begin{center}
\begin{minipage}{\textwidth}
\makeatletter
\def\@captype{figure}
\makeatother
\centering
\[
\begin{tikzcd}[column sep=large, row sep=large]
\{\text{Weil 除子}\;D\}
 \arrow[r, leftrightarrow, "D\mapsto I_D"]
 \arrow[d, two heads, "\text{模主除子}"'] &
\{\text{反射分式理想}\;I\subset K\}
 \arrow[d, two heads, "\text{模}\;K^*\;\text{倍乘}"] \\
\operatorname{Cl}(R) \arrow[r, "\sim"] &
\{\text{秩一反射模同构类}\} \\
\operatorname{Pic}(R) \arrow[u, hook] \arrow[r, "\sim"] &
\{\text{逐点自由秩一的模同构类}\}
 \arrow[u, hook]
\end{tikzcd}
\]
\[
\begin{aligned}
\operatorname{Cl}(R)=0
&\quad\Longleftrightarrow\quad R\;\text{为 UFD},\\
\text{本节单射}\;\operatorname{Pic}(R)\hookrightarrow\operatorname{Cl}(R)\;\text{为同构}
&\quad\Longleftrightarrow\quad R\;\text{局部阶乘}.
\end{aligned}
\]
\caption[除子、反射模与 Picard 群]{除子、反射模与 Picard 群。\(R\) 为交换 Noether 正规整环，\(K\) 为其分式域，图中的模均有限生成，分式理想均非零。箭头采用本节的理想正号约定。}
\label{b3:fig:divisor-reflexive-picard}
\end{minipage}
\end{center}

图中的两条向下满射分别对主除子和非零标量倍乘取商，因此中间一行是除子类与模同构类的对应。最后一行只收取逐点自由秩一的模类；它们嵌入所有秩一反射模类中，局部阶乘性恰好使这两个集合相等。


\begin{chaptersummary}
Serre 判据中，\(R_1\) 提供高度一处的 DVR，\(S_2\) 保证一个分式不属于环的障碍必在这些素点处显现。对有多个不可约分量的环，还须借助全分式环中的幂等元，证明每个正规局部环只有一个极小素理想。因此非零 Noether 正规环是有限个整闭整环的直积。

\ProseParagraphBreak
正规整环是其高度一局部环在分式域中的交；各处整数阶组成 Weil 除子\footnote{韦伊（André Weil，1906-1998），法国数学家，研究代数几何与数论，奠定抽象代数几何的基础。}，主除子记录同一个有理函数的全部阶数。类群为零恰好等价于 UFD。秩一反射模保存这些高度一数据，而可逆模还要求在全部素点处局部自由。正规锥面的非零类及 Dedekind 环的可逆非主理想分别显示，这些局部与整体条件确有差别。
\end{chaptersummary}

\BookExercises

\begin{exercise}\label{b3:ex:20-r0-s1-tests}
设 \(k\) 为域，\(A=k[x,y]/(x^2,xy^2)\)。求 \(A\) 的幂零根、该幂零根的支集及 \(\operatorname{Ass}_AA\)，判断 \(R_0\)、\(S_1\) 是否成立，并计算幂零根在其支集处的长度。
\end{exercise}
\begin{solution}
幂零根为 \(N=(\bar x)\)，且 \(\operatorname{Ann}_AN=(\bar x,\bar y^2)\)，故 \(\operatorname{Supp}_AN=\{\mathfrak m\}\)，其中 \(\mathfrak m=(\bar x,\bar y)\)。\(N\) 作为 \(k\) 向量空间以 \(\bar x,\bar x\bar y\) 为基，在 \(\mathfrak m\) 处的长度为二。\(\bar x\bar y\ne0\) 的零化子为 \(\mathfrak m\)，所以 \(\mathfrak m\) 为关联素理想。正合列 \(0\to N\to A\to k[y]\to0\) 则使所有关联素理想包含在 \(\{\mathfrak m,(\bar x)\}\) 中；唯一极小素理想 \((\bar x)\) 必为关联素理想，故恰有这两个。在 \((\bar x)\) 处，\(y\) 可逆而 \(xy^2=0\)，所以 \(\bar x=0\)，局部化为 \(k(y)\)，\(R_0\) 成立。嵌入关联素理想 \(\mathfrak m\) 使 \(S_1\) 失败。这里未约化部分集中于闭点，但其长度为二，不能只凭一般点为域辨认它。
\end{solution}

\begin{exercise}\label{b3:ex:20-one-dimensional-s2}
证明每个一维交换 Noether 整环满足 \(S_2\)。由此解释为什么在一维中，正规性仍需额外检查 \(R_1\)。
\end{exercise}
\begin{solution}
高度零处的局部环为分式域；高度一处为一维局部整环，可从极大理想取非零元素作为非零因子，因此深度至少一。这恰满足 \(\operatorname{depth}R_{\mathfrak p}\ge\min(2,\dim R_{\mathfrak p})\)。尖点整环满足这个条件，但高度一特殊点的极大理想需要两个生成元，不是 DVR，故不正规。
\end{solution}

\begin{exercise}\label{b3:ex:20-normal-product}
设 \(A,B\) 为非零交换 Noether 整闭整环。证明 \(A\times B\) 正规而不是整环，并说明其两个不可约分量为何不相交。
\end{exercise}
\begin{solution}
两个非零幂等元 \((1,0),(0,1)\) 的乘积为零，故不是整环。任一素理想必包含其中恰一个：乘积为零迫使至少包含一个，和为一又排除同时包含。于是局部环来自 \(A\) 或 \(B\)，并且整闭，所以直积正规。两个不可约分量分别是 \(V(A\times0)\) 与 \(V(0\times B)\)，它们的交为 \(V(A\times B)=\varnothing\)，且各自都是开闭子集。
\end{solution}

\begin{exercise}\label{b3:ex:20-normal-dimension-one}
设 \(R\) 为维数至多一的交换 Noether 环。证明 \(R\) 正规当且仅当所有局部环都正则。
\end{exercise}
\begin{solution}
正规时，每个零维局部环为域，每个一维局部环为 DVR，所以正则。反向，正则局部环为 UFD，故整闭整环，于是按定义正规。也可以在维数至多一时用正则性提供 \(R_1\) 和 CM 性提供 \(S_2\)，再应用 Serre 判据。
\end{solution}

\begin{exercise}\label{b3:ex:20-principal-local-test}
设 \(R\) 为交换 Noether 正规整环，\(0\ne a,b\in R\)，对每个高度一素理想 \(\mathfrak p\)，记 \(v_{\mathfrak p}\) 为 \(R_{\mathfrak p}\) 的归一化离散赋值。证明 \(a\mid b\) 当且仅当每个高度一 \(\mathfrak p\) 均满足 \(v_{\mathfrak p}(a)\le v_{\mathfrak p}(b)\)。由此给出两个非零元素相伴的赋值判据。
\end{exercise}
\begin{solution}
整除性等价于 \(b/a\in R\)，再用高度一交定理即可。若两元素在所有高度一处的阶数相等，则 \(b/a\) 与 \(a/b\) 均属于 \(R\)，因而 \(b/a\) 为单位，两元素相伴；反向单位在每个赋值处阶数为零。
\end{solution}

\begin{exercise}\label{b3:ex:20-fraction-divisor}
设 \(k\) 为域。在 \(R=k[x,y]\) 中计算
\(f=x^2y^3/(x-y)^4\) 的主除子，并求使 \(f\in R_{\mathfrak p}\) 的高度一素理想。
\end{exercise}
\begin{solution}
\(x,y,x-y\) 是两两不相伴的素元，所以
\(\operatorname{div}(f)=2[(x)]+3[(y)]-4[(x-y)]\)。在 \((x-y)\) 处阶数为负，其他所有高度一素理想处阶数非负，因此恰在 \(\mathfrak p\ne(x-y)\) 时属于该局部环。这些结论不受域特征影响。
\end{solution}

\begin{exercise}\label{b3:ex:20-remove-codimension-two}
设 \(R\) 为 Noether 正规整环，\(Z\subseteq\operatorname{Spec}R\) 为不含任何高度一素点的闭集。若 \(f\in\operatorname{Frac}(R)\) 满足 \(f\in R_{\mathfrak p}\) 对所有 \(\mathfrak p\notin Z\)，证明 \(f\in R\)。说明这里并未断言每个理想都由相同的局部数据恢复。
\end{exercise}
\begin{solution}
由于 \(\operatorname{Spec}R\setminus Z\) 包含所有高度一素点，题设使 \(f\in R_{\mathfrak p}\) 对每个高度一 \(\mathfrak p\) 成立。由\cref{b3:thm:normal-height-one-intersection}，\(f\in R\)。对于理想则不成立：对域 \(k\)，在 \(k[x,y]\) 中，\((x,y)\) 在闭点 \((x,y)\) 之外均变为整个局部环，但它本身是真理想。函数的可延拓性与任意子模的可恢复性不是同一个命题。
\end{solution}

\begin{exercise}\label{b3:ex:20-localization-class-group}
设 \(R\) 为交换 Noether 正规整环，\(P\) 为高度一素理想。证明可以选取 \(0\ne a\in P\)，使 \(v_P(a)=1\)，并由 \(\operatorname{div}(a)\) 把 \([P]\) 表为其他高度一素理想类的整数线性组合。再用局部化到 \(R_P\) 的类群映射解释这一现象。
\end{exercise}
\begin{solution}
\(R_P\) 的统一化元可写成 \(a/s\)，其中 \(a\in R\)、\(s\notin P\)。\(s\) 在此处的赋值为零，所以 \(v_P(a)=1\)，特别地 \(a\in P\)。\(a\) 是环元素，在所有高度一素点的赋值均非负，且只有有限多个非零。因此可写
\[
\operatorname{div}(a)=[P]+\sum_{Q\ne P}n_Q[Q],\qquad n_Q\ge0,
\]
其中求和仅含有限多个高度一素理想。主除子在类群中为零，故 \([P]=-\sum_{Q\ne P}n_Q[Q]\)。取乘法集 \(S=R\setminus P\)，\(R_P\) 是 DVR，类群为零；\cref{b3:prop:class-group-localization}说明整个 \(\operatorname{Cl}(R)\) 由与 \(S\) 相交的高度一素理想类生成。不同的高度一素理想互不包含，所以这些恰好是 \(Q\ne P\)。上面的主除子计算给出了 \([P]\) 如何由它们表示的具体关系。
\end{solution}

\begin{exercise}\label{b3:ex:20-quadric-three-divisors}
设 \(k\) 为域。在二维正规锥面局部环 \(A=k[x,y,z]_{(x,y,z)}/(xy-z^2)\) 中，记 \(P=(x,z)\)、\(Q=(y,z)\)。计算 \(\operatorname{div}(x)\)、\(\operatorname{div}(y)\)、\(\operatorname{div}(z)\)，并比较 \([P]\) 与 \([Q]\)。
\end{exercise}
\begin{solution}
含 \(x\) 的唯一高度一素理想为 \(P\)，在其处 \(y\) 为单位、\(x=z^2/y\)，故 \(\operatorname{div}(x)=2[P]\)。对称地 \(\operatorname{div}(y)=2[Q]\)。含 \(z\) 的高度一素理想由 \(xy=0\) 得到 \(P,Q\)，且 \(z\) 在两处都是统一化元，故 \(\operatorname{div}(z)=[P]+[Q]\)。因此两类互为逆，而各自二倍为零，所以 \([P]=[Q]\)。正文证明它们不是零类。
\end{solution}

\begin{exercise}\label{b3:ex:20-regular-quotient-loses-s1}
取\cref{b3:exm:r1-not-s2}的两平面局部环 \(R\)，令 \(a=x+z\)。证明 \(a\) 为非零因子，但 \(R/aR\) 不满足 \(S_1\)。
\end{exercise}
\begin{solution}
\(R\) 约化，其两个极小素理想分别对应两张平面。\(x+z\) 在这两个素理想的商环中的像分别为 \(x\) 或 \(z\)，均非零，所以避开全部关联素理想，为非零因子。正文已得 \(\dim R=2\)、\(\operatorname{depth}R=1\)，由\cref{b3:prop:regular-quotient-dimension,b3:cor:regular-quotient-depth}，商的维数为一、深度为零。于是其闭点是关联素理想却不是极小素理想，\(S_1\) 失败。
\end{solution}

\begin{exercise}\label{b3:ex:20-height-one-ideal-cm}
设 \((R,\mathfrak m)\) 为二维 Noether 正规局部整环，\(P\) 为高度一素理想。证明 \(P\) 作为 \(R\) 模为极大 CM 模。
\end{exercise}
\begin{solution}
Serre 判据给出 \(\operatorname{depth}R=2\)。\(R/P\) 为一维局部整环，深度为一；维数等于一是因为 \((0)\subsetneq P\subsetneq\mathfrak m\) 且 \(\dim R=2\)。正合列 \(0\to P\to R\to R/P\to0\) 的深度引理给出 \(\operatorname{depth}P\ge\min(2,1+1)=2\)。非零理想 \(P\) 的零化子为零，维数二，所以为极大 CM 模。
\end{solution}

\begin{exercise}\label{b3:ex:20-normal-birational-finite}
设 \(R\subseteq B\subseteq K=\operatorname{Frac}(R)\)，其中 \(R\) 为 Noether 正规整环，\(B\) 为有限生成的 \(R\) 模的子环。证明 \(B=R\)。解释为什么这不能推出正规空间没有奇点。
\end{exercise}
\begin{solution}
有限代数的每个元素对 \(R\) 整，且全部位于 \(K\)，整闭性使 \(B\subseteq R\)，于是相等。这只排除了同一分式域中的有限整扩张；二维正规锥面仍有非正则闭点。消解奇点使用的几何改变不必由一个这样的有限子环扩张给出，因而不受这个结论排除。
\end{solution}

\begin{exercise}\label{b3:ex:20-cusp-normality-check}
设 \(k\) 为域，令 \(R=k[t^3,t^4,t^5]_{(t^3,t^4,t^5)}\)。证明它满足 \(S_2\) 但不正规，并给出一个整而不属于 \(R\) 的分式。
\end{exercise}
\begin{solution}
该环有限整扩张于 \(k[t^3]\) 的相应局部环，维数一；它是整环，故满足 \(S_2\)。\(t=t^4/t^3\) 位于分式域，且满足 \(T^3-t^3=0\)。若 \(t=a/b\) 且 \(b(0)\ne0\)，则 \(tb\) 有非零一次项，而半群 \(\langle3,4,5\rangle\) 中没有一，所以 \(a\) 没有一次项，矛盾。故不正规。
\end{solution}

\begin{optionalexercise}\label{b3:ex:20-reflexive-maximal-ideal}
设 \(k\) 为域，在 \(R=k[x,y]\) 中令 \(I=(x^3,x^2y,xy^2,y^4)\)。求 \(I^{**}\)，计算 \(\dim_k(R/I)\)，并判断 \((I^n)^{**}\) 是否随正整数 \(n\) 改变。说明这些双对偶保留了哪些局部数据。
\end{optionalexercise}
\begin{solution}
每个高度一素理想都不能同时包含 \(x,y\)，因此包含在 \(I\) 中的 \(x^3,y^4\) 至少有一个在该处可逆，\(I_{\mathfrak p}=R_{\mathfrak p}\)。高度一交公式给出 \(I^{**}=R\)。商 \(R/I\) 的单项式基为
\[
1,x,x^2,y,xy,y^2,y^3,
\]
所以维数为七，求值映射的余核支集为闭点 \((x,y)\)。对每个 \(n\ge1\)，\(I^n\) 包含 \(x^{3n},y^{4n}\)，故它在全部高度一素点处仍为单位理想，\((I^n)^{**}=R\)。这些理想在高度二闭点处的消失条件不同，双对偶却都只恢复共同的高度一局部数据。
\end{solution}

\begin{optionalexercise}\label{b3:ex:20-reflexive-product-quadric}
设 \(k\) 为域，在锥面局部环 \(A=k[x,y,z]_{(x,y,z)}/(xy-z^2)\) 中令 \(P=(x,z)\)、\(\mathfrak m=(x,y,z)A\)。对每个整数 \(r\ge1\)，分别计算 \(P^{2r}\)、\(P^{2r+1}\) 及其双对偶，并证明这两种普通幂均严格小于相应的双对偶。
\end{optionalexercise}
\begin{solution}
由正文的 \(P^2=x\mathfrak m\)，得到
\[
P^{2r}=x^r\mathfrak m^r,\qquad
P^{2r+1}=x^r\mathfrak m^rP.
\]
在每个高度一素点，\(\mathfrak m\) 变为单位理想。因此两式的高度一局部数据分别与 \(x^rA\)、\(x^rP\) 相同。\(P\) 为反射理想，非零标量倍乘保持反射性，故
\[
(P^{2r})^{**}=x^rA,\qquad
(P^{2r+1})^{**}=x^rP.
\]
第一项严格包含来自 \(\mathfrak m^r\subsetneq A\)。若第二项相等，整环中约去 \(x^r\) 得 \(P=\mathfrak m^rP\subseteq\mathfrak mP\)，Nakayama 引理将给出 \(P=0\)，矛盾。因而普通幂一直保留闭点处的额外消失条件；反射化后的幂则按高度一阶数的奇偶分别为主理想或 \(P\) 的标量倍。
\end{solution}

\begin{optionalexercise}\label{b3:ex:20-rank-one-dual}
设 \(k\) 为域，在锥面局部环 \(A=k[x,y,z]_{(x,y,z)}/(xy-z^2)\) 中令 \(P=(x,z)\)、\(\mathfrak m=(x,y,z)A\)。在分式域中计算 \(P^*=(A:P)\)，并求求值配对 \(P\otimes_AP^*\to A\) 的像与余核。说明为什么 \((P\otimes_AP^*)^{**}\cong A\) 不要求原配对满射。
\end{optionalexercise}
\begin{solution}
\(P\) 只在高度一素理想 \(P\) 处具有阶数一，其他高度一处阶数为零；\(x\) 的主除子为 \(2[P]\)。因此 \(x^{-1}P\) 在这些素点的阶数恰为 \(P^*\) 所需的相反数。两者均为反射分式理想，高度一交公式给出
\[
P^*=x^{-1}P=A+A\cdot\frac zx.
\]
也可直接检查：\(z/x\) 乘 \(x,z\) 分别得到 \(z,y\)，确实把 \(P\) 送入 \(A\)。求值配对在此表示下是乘法，其像为
\[
PP^*=(x,z)\left(1,\frac zx\right)=(x,z,y)=\mathfrak m,
\]
余核为剩余域 \(k\)，故并不满射。另一方面，\(\mathfrak m\) 在每个高度一素点处变为单位理想，\(\mathfrak m^{**}=A\)。乘法 \(P\otimes_AP^*\to\mathfrak m\) 的核为挠子模，取双对偶先零化这部分，再补回闭点缺少的元素，得到 \((P\otimes_AP^*)^{**}\cong A\)。
\end{solution}

\begin{optionalexercise}\label{b3:ex:20-reflexive-localization}
设 \(R\) 为 Noether 正规整环，\(M\) 为有限生成反射模，\(S\subseteq R\setminus\{0\}\) 为乘法集。证明 \(S^{-1}M\) 为 \(S^{-1}R\) 上的反射模。
\end{optionalexercise}
\begin{solution}
Noether 性保证 \(M\) 与 \(M^*\) 均有限展示，故两次使用 Hom 的局部化同构得到
\(S^{-1}(M^{**})\cong(S^{-1}M)^{**}\)。在这个同构下，原求值映射的局部化正是新的求值映射；前者由反射性为同构，因此后者也为同构。
\end{solution}

\begin{optionalexercise}\label{b3:ex:20-picard-versus-class}
给出三种例子并说明理由：局部阶乘但不是 UFD；UFD 但不是主理想整环；正规局部环且 Picard 群为零，但除子类群非零。
\end{optionalexercise}
\begin{solution}
第一种取 \(\mathbb Z[\sqrt{-5}]\)，它是 Dedekind 整环，故各非零素点为 DVR；第10章的可逆非主理想 \((2,1+\sqrt{-5})\) 使类群非零，故不是 UFD。第二种取 \(k[x,y]\)，它为 UFD，而 \((x,y)\) 不主生成。第三种取正规锥面局部环；局部环上的秩一有限投射模自由，Picard 群为零，但 \((x,z)\) 给出阶二的非零除子类。
\end{solution}
